\documentclass[11pt,a4paper]{article}
\usepackage[T1]{fontenc}
\usepackage[utf8]{inputenc}
\usepackage{lmodern}
\usepackage[a4paper,margin=25mm,headheight=15pt]{geometry}
\usepackage{amsmath,amssymb,amsthm,mathtools}
\usepackage{microtype,booktabs,longtable,array,xcolor,listings,fancyhdr}
\usepackage[hyphens]{url}
\usepackage{xurl}
\usepackage[colorlinks=true,linkcolor=blue!45!black,citecolor=blue!45!black,urlcolor=blue!45!black]{hyperref}
\hypersetup{pdftitle={A smaller iPRA presentation by first-instance witness discharge},pdfauthor={Research report prepared for Vladimir Reshetnikov}}
\newtheorem{theorem}{Theorem}[section]
\newtheorem{lemma}[theorem]{Lemma}
\newtheorem{proposition}[theorem]{Proposition}
\newtheorem{corollary}[theorem]{Corollary}
\theoremstyle{definition}\newtheorem{definition}[theorem]{Definition}
\theoremstyle{remark}\newtheorem{remark}[theorem]{Remark}
\newcommand{\ipra}{\mathsf{iPRA}}
\newcommand{\nf}{\operatorname{nf}}
\newcommand{\FV}{\operatorname{FV}}
\pagestyle{fancy}\fancyhf{}
\fancyhead[L]{\small First-instance witness discharge for iPRA}
\fancyhead[R]{\small\thepage}
\renewcommand{\headrulewidth}{0.3pt}
\setlength{\emergencystretch}{2em}
\allowdisplaybreaks[2]
\raggedbottom
\lstset{language=Python,basicstyle=\footnotesize\ttfamily,columns=fullflexible,breaklines=true,breakatwhitespace=false,showstringspaces=false,frame=single,rulecolor=\color{black!20},tabsize=4}
\title{A smaller presentation of intuitionistic\newline primitive recursive arithmetic\\[4pt]
\large First-instance witness discharge}
\author{Research report prepared for Vladimir Reshetnikov}
\date{11 November 2026}
\begin{document}
\maketitle
\begin{abstract}
The current witness-discharge presentation of full first-order iPRA has
three inference families. Its elimination rule takes a major formula,
a local equality proof $b=t$, and a local formula proof $F$, then exports
the normalized instance $F[t/b]$. We prove that the equality proof is
unnecessary as a primitive premise: a deterministic test of the opened
clause's zero and first instances supplies the only substitution that
the rule needs to perform. The new rule has two proof premises and no
replacement-term argument. Nevertheless it preserves every open sequent,
not merely the closed equations.

The reconstruction obtains ordinary elimination and numerical cases
before equality, and then recovers arbitrary transport through a
two-row naming matrix. The inherited full-iPRA bi-interpretation therefore
has both numerical domains whole, actual equality, and identity numerical
composites. We include the complete primitive specification, proofs,
the normalization and arithmetic bridges, a standalone Python checker,
proof translations, adversarial examples, and reproducible measurements.
The result is a concrete simplification relative to the current
witness-discharge presentation, not an absolute minimality theorem.
\end{abstract}
\tableofcontents
\clearpage

\section{Result and comparison boundary}
\label{sec:boundary}

The comparison repository is ProveIt~\cite{baseline}, under the subtree
\begin{center}\small\path{Logic/PrimitiveRecursiveArithmetic/ConciseIPRA}.\end{center}
The pinned commit is
\begin{center}\small\ttfamily 17898ea7e6a1300dfdcad71040d5daeed67916d4.\end{center}
Its README lists sixty-seven filed research packages. We also inspected
the six archives in \path{docs/incoming}, including the recent
\texttt{witness\_discharge} and \texttt{uniform\_binding} packages.
The immediate mathematical control is the incoming
\path{ipra_witness_discharge_2026-11-11.zip}~\cite{witness}.
In particular, three inference families and an initial reflexivity clause
are already present there. The new theorem concerns deletion of one
\emph{proof premise} and deterministic recovery of its substitution.

The complete proposed calculus is denoted $\mathcal A$.
Its formulas have only equality and the paired binder $D$; its numerical
terms are canonical scalar iterators. There is one initial clause and
three inference families: Use, Intro, and Name. The standard reading is
\[
  D_{x,y}B(x,y)
    \quad\rightsquigarrow\quad
  \exists x\bigl(B(x,0)\to\forall y\,B(x,y)\bigr).
\]
The syntactic reductions used to obtain canonical terms are explicitly
part of the theory's specification and of the measured checker.

\paragraph{Reference theory.}
Here $\ipra$ means full intuitionistic first-order logic with equality,
zero, successor, and a function symbol for every fixed primitive-recursive
description, the usual defining equations and successor axioms, and
induction restricted to quantifier-free formulas, with parameters.
It is not the quantifier-free proof language alone. The equivalent
atomic-induction convention in the full primitive-recursive language is
also used in the literature~\cite{khaniki}; the equality-reflection
appendix gives the required equivalence in this development.

\paragraph{What is proved.}
For every canonical open clause sequent, the old witness-discharge
calculus $\mathcal W$ and $\mathcal A$ prove exactly the same conclusion
from exactly the same assumptions. The translation is the identity on
their numerical and formula languages. Composing with the scalar
arithmetic bridge proves full bi-interpretability and equiconsistency
with $\ipra$. The two numerical composite interpretations are internally
isomorphic through the graph $x=y$; both formula composites are also
proved, so this is stronger than a standard-model agreement or an
equational conservativity statement.

\paragraph{What is inherited.}
The scalar iterator, the $D$ binder, the first-order logical macros,
full computation normalization, the two-row naming method, predecessor,
equality reflection, and the primitive-recursive compiler are inherited
from the repository. The new contribution is the completeness of a
fixed first-instance restriction on witness discharge. The transparent
certificate representation and its complete implementation are also
given, with matched controls. No universal priority or absolute
minimum-source-size claim is made.

\paragraph{How to read the report.}
Section~\ref{sec:deterministic-name} contains the complete native
axiomatization and the new proof. Section~\ref{sec:canonical-bridge}
connects canonical syntax to the raw scalar basis. The implementation
section records the complete checker contract and its measured scope.
The appendices reproduce the attributed logical and arithmetic
foundations, so the full-iPRA conclusion does not rest on an omitted
predecessor or pairing argument.


% Mathematical fragment for deterministic clause-derived witness discharge.
% Prepared in this session by the logical-reduction subagent.
% Inherited sources: incoming ipra_witness_discharge_2026-11-11,
% especially report.tex sections 2.1--2.8 and validation/macros.py.
% The new result is the deterministic two-premise rule, not the scalar bridge.

\section{Discharging a witness using its first clause instance}
\label{sec:deterministic-name}

The incoming witness-discharge presentation has a separate proof of an
equality $b=t$ inside each eliminated witness scope. That premise is
unnecessary. It suffices to recognize an equality in one fixed instance of
the clause which the witness itself supplies. If that instance has no usable
equality, the witness is discharged without substitution. The choice is
determined by the major formula; the certificate supplies neither an
equality proof nor a replacement term.

This is a restriction of the previous inference, with a nontrivial
completeness theorem. It does not rely on the numerical standard model,
semantic truth testing, or an unlisted equality rule.

\subsection{Complete native specification}

Use one numerical sort. Let $N$ denote full computation normalization of
the scalar terms
\[
 t ::= x\mid 0\mid St\mid I(n,a;u.h)
\]
for the context-closed reductions
\[
 I(0,a;u.h)\longrightarrow a,
 \qquad I(Sn,a;u.h)\longrightarrow h[I(n,a;u.h)/u].
 \tag{NF}
\]
Only the state variable $u$ in $h$ is bound. Every displayed term means
its normal form, and every substitution is capture-avoiding and followed
by $N$. Normalization of a formula acts on numerical terms and leaves its
formula constructors intact. We use the established termination,
confluence, and substitution-coherence results for these reductions.
The normalization algorithm is part of the trusted syntax and must be
included in a complete-checker measurement.

Formulas are
\[
 A ::= (s=t)\mid D_{x,y}A,
\]
where $x,y$ are distinct bound numerical variables. A reusable context
clause is $[\vec A;y.B]$, where $\vec A$ has zero or one antecedent,
and $y$ is absent from every antecedent. Its ordinary reading is
\[
 [\,;y.B]^\circ=\forall y B^\circ,
 \qquad [A;y.B]^\circ=A^\circ\to\forall y B^\circ.
\]
An ordinary assumption $A$ abbreviates $[\,;v.A]$, with $v$ unused.
Contexts are finite lists of clause occurrences. Every proof starts with
the single explicit initial clause
\[
                 \mathcal R=[\,;y.y=y]. \tag{Ax}
\]
Put $H=(0=0)$, $1=S0$, and $2=S1$. The designated contradiction is $0=1$.
There are exactly the following three primitive inference families.

\paragraph{Use.}
For a retained $[\vec A;y.B]$ and any well-scoped term $t$, infer $N(B[t/y])$
after proving the antecedent if there is one. There is no proof premise
for a clause with no antecedent.
\[
 \frac{\{\Gamma\vdash A:A\text{ is the sole antecedent, if present}\}}
      {\Gamma\vdash N(B[t/y])}.
 \tag{Use}
\]

\paragraph{Intro.}
For fresh $z$ and a well-scoped witness $t$,
\[
 \frac{\Gamma,N(B(t,\eta))\vdash N(B(t,Sz))}
      {\Gamma\vdash D_{x,y}B(x,y)},
 \qquad
 \eta=\begin{cases}
 z,& B\text{ is an equation and }t\text{ is literally }0,\\
 0,&\text{otherwise}.
 \end{cases}
 \tag{Intro}
\]
Here $z$ is fresh for $\Gamma,t$ and all fixed parameters of $B$.
The displayed ordinary assumption is discharged.

\paragraph{Name.}
After proving $D_{x,y}B$, take a fresh numerical parameter $b$ and form
\[
 B_b(y)=N(B(b,y)),\qquad
 A_0=N(B_b(0)),\qquad A_1=N(B_b(1)),
 \qquad \mathcal B_b=[A_0;y.B_b(y)].
\]
Define the replacement term \emph{deterministically} by
\[
 \tau_b(B)=
 \begin{cases}
 t,& A_0\text{ is literally }H\text{ and }A_1
       \text{ is literally the equation }b=t,\\
 b,&\text{otherwise}.
 \end{cases}
 \tag{First-instance test}\label{eq:first-instance-test}
\]
The literal equation test is ordered: a formula $t=b$ is not recognized
unless its left side is itself literally $b$. The term $t$ is the right
side of $A_1$, so no separate term argument is supplied. Then the rule is
\[
 \frac{\Gamma\vdash D_{x,y}B(x,y)\qquad
       \Gamma,\mathcal B_b\vdash F}
      {\Gamma\vdash N(F[\tau_b(B)/b])},
 \tag{Name}\label{eq:deterministic-name}
\]
provided the exported normal form has no free $b$. The parameter $b$ is
fresh for $\Gamma$ and every fixed parameter of $B$. The branch formula
$F$ is well scoped in the extended numerical context. A recognized $t$
may contain $b$; substitution is simultaneous, so an inserted copy of $t$
is not substituted again.

These are all the primitive inferences. Reflexivity is a Use of
$\mathcal R$. Computation is accounted for by canonical syntax. There
are no additional primitive equality, logical, induction, case, cut, or
weakening rules.

\subsection{Soundness and the escape condition}

The ordinary formula reading is
\[
 (s=t)^\circ=(s=t),\qquad
 (D_{x,y}B)^\circ=
       \exists x\bigl(B^\circ(x,0)\to\forall y B^\circ(x,y)\bigr).
\]
For this subsection, use intuitionistic first-order equality logic with
the scalar computation equations, numerical coverage, separation
$0=1\to\bot$, and contextual induction for literal equations.
This is the inherited scalar theory with a full-iPRA bi-interpretation.

\begin{lemma}[The replacement equality is always available]
\label{lem:det-replacement}
In the context $\Gamma,\mathcal B_b$, one can derive
$b=\tau_b(B)$ using only the initial reflexivity clause and Use.
\end{lemma}
\begin{proof}
If the test fails, $\tau_b(B)=b$ and Use of $\mathcal R$ at $b$ proves
$b=b$. If it succeeds, the antecedent of $\mathcal B_b$ is exactly
$H$. Use of $\mathcal R$ at zero proves $H$, and Use of $\mathcal B_b$
at $1$ proves its normalized first instance, which is exactly
$b=\tau_b(B)$. No equality reasoning occurs in either derivation.
\end{proof}

\begin{proposition}[Soundness of Name]
Rule~\eqref{eq:deterministic-name} is sound for the ordinary reading,
including when the recognized replacement term contains $b$.
\end{proposition}
\begin{proof}
Ordinary existential elimination opens the major formula at $b$ with
the assumption represented by $\mathcal B_b$. The branch gives $F$.
Lemma~\ref{lem:det-replacement} gives $b=\tau_b(B)$, and ordinary equality
substitution gives $F[\tau_b(B)/b]$. Computation equations identify this
with the exported normal form. That normal form contains no free $b$,
so ordinary existential elimination discharges the witness and its clause.
Equality substitution itself does not require $b$ to be absent from its
right side; only existential discharge needs the stipulated absence from
the final exported formula.
\end{proof}

Use has its displayed ordinary reading. Intro is sound by numerical
coverage in case mode and literal-equation induction in progress mode.
Thus the entire presentation is sound in the inherited scalar theory.

\subsection{Ordinary elimination and numerical cases without equality}

\begin{lemma}[Ordinary paired elimination]
\label{lem:det-ordinary-elim}
The usual rule
\[
 \frac{\Gamma\vdash D_{x,y}B\qquad
       \Gamma,\mathcal B_b\vdash A}
      {\Gamma\vdash A},\qquad b\notin\operatorname{FV}(A),
\]
is derivable with one use of Name.
\end{lemma}
\begin{proof}
Apply Name with its supplied branch formula $A$. Since $A$ is a normal
formula with no free $b$, substituting any term for $b$ leaves it unchanged.
Thus $N(A[\tau_b(B)/b])=A$ regardless of whether the first-instance test
succeeds. There is no need to force the test to fail or to choose a target.
\end{proof}

Weakening, injective fresh renaming, and ordinary-assumption cut are
admissible by induction on proof trees. Freshen local eigenparameters
before inserting proofs. The deterministic test depends only on the
major formula and is unchanged by an injective renaming. An ordinary
assumption is an unconditional constant clause, so its Use nodes can be
replaced by the supplied proof. None of these structural constructions
uses general equality transport or arbitrary numerical substitution of
proof trees.

For unused $u$, write $q_y A(y)=D_{u,y}A(y)$.
\begin{lemma}[Conditional closure and numerical cases]
\label{lem:det-cases}
The calculus derives
\[
 \frac{\Gamma,A(0)\vdash A(Sz)}{\Gamma\vdash q_yA(y)},
 \qquad
 \frac{\Gamma\vdash q_yA(y)\qquad\Gamma\vdash A(0)}
      {\Gamma\vdash A(t)},
\]
and consequently
\[
 \frac{\Gamma\vdash A(0)\qquad\Gamma\vdash A(Sz)}
      {\Gamma\vdash A(t)}
 \tag{Cases}\label{eq:det-cases}
\]
for every native formula $A$ and appropriate fresh $z$.
These constructions use no equality transport or equation induction.
\end{lemma}
\begin{proof}
For the first rule, introduce the unused witness at the numeral $1$.
It is not literally zero, so Intro uses the case premise at every matrix
root. For the second, use Lemma~\ref{lem:det-ordinary-elim}. The opened
clause is $[A(0);y.A(y)]$, because the witness coordinate is unused.
Supply its antecedent and apply Use at $t$. The result has no free
existential witness, so ordinary elimination closes. For Cases, first
weaken the successor proof by the unused ordinary assumption $A(0)$,
then combine the preceding two constructions.
\end{proof}

\subsection{Recovering equality from the fixed first instance}

Define normalizing selector and parity terms by
\[
 C(v,a,c)=N(I(v,a;u.c)),\qquad u\notin\operatorname{FV}(c),
 \qquad P(y)=N(I(y,0;u.C(u,1,0))).
\]
Direct computation gives literal identities of canonical syntax:
\[
 C(0,a,c)\equiv a,\qquad C(Sv,a,c)\equiv c,
 \qquad P(1)\equiv1,\qquad P(2)\equiv0.
 \tag{Endpoints}\label{eq:det-endpoints}
\]
Only these fixed endpoints of $P$ are used; no parity induction or range
theorem is required.

\begin{theorem}[Uncrossed term transport]
\label{thm:det-transport}
For every fixed numerical template $h(w)$, the rule
\[
 \frac{\Gamma\vdash s=t\qquad\Gamma\vdash N(h(s))=r}
      {\Gamma\vdash N(h(t))=r}
\]
is derivable. It uses no new arithmetic assumption and no supplied
equality branch at Name.
\end{theorem}
\begin{proof}
Choose $x,y$ fresh for all passive data. Let $h(x)$ denote its canonical
instance and form the equation matrix
\[
 \begin{split}
 B(x,y)=\bigl(&C(y,0,C(P(y),N(h(x)),x))\\
             &{}=C(y,0,C(P(y),r,t))\bigr).
 \end{split}
 \tag{Naming matrix}\label{eq:det-naming-matrix}
\]
Introduce $D_{x,y}B(x,y)$ at witness $s$. Its successor instance is
\[
 C(P(Sz),N(h(s)),s)=C(P(Sz),r,t).
\]
Apply Cases to the template
\[
       A(v)=\bigl(C(v,N(h(s)),s)=C(v,r,t)\bigr)
\]
at target $P(Sz)$. Its zero instance is exactly the supplied equation
$N(h(s))=r$, and its successor instance is exactly the supplied equation
$s=t$. The eigenparameter used for this Cases argument is fresh for the
current local context and is different from $z$.

The resulting successor proof ignores the ordinary assumption of Intro.
Consequently it is valid whether $s$ is literally zero, selecting progress,
or any other term, selecting cases. This constructs the major formula
without any equality substitution.

Open that major at fresh $b$. By the endpoint computations,
\[
 B(b,0)\equiv H,\qquad B(b,1)\equiv(b=t),
 \qquad B(b,2)\equiv(N(h(b))=r).
\]
The first-instance test therefore selects exactly $\tau_b(B)=t$.
For the sole branch of Name, prove $H$ by the fixed reflexivity clause
and specialize the opened clause at $2$. This yields
$F=(N(h(b))=r)$. Name exports
\[
 N\bigl((N(h(b))=r)[t/b]\bigr)
       \equiv (N(h(t))=r)
\]
by normalization--substitution coherence. The parameter $b$ was fresh
for $h,t,r$, so the exported formula is outer scoped.

The premise carrying $s=t$ is used while proving the major naming formula,
not as a third premise of Name. Thus the proof does not smuggle the deleted
equality branch back into the rule's checking interface.
\end{proof}

\begin{corollary}[All equality reasoning]
\label{cor:det-equality}
Symmetry, transitivity, congruence, crossed transport, and transport in
every native formula template are derivable.
\end{corollary}
\begin{proof}
For symmetry, take $h(w)=w$ and $r=s$ in
Theorem~\ref{thm:det-transport}; the minor premise is reflexivity $s=s$.
This proves $t=s$ from $s=t$. Uncrossed transport followed by symmetry
gives crossed transport. A reflexive minor gives congruence, followed by
symmetry if necessary. Symmetry of $a=b$ and transport of $b=c$ through
the identity template give $a=c$.

For transport in an equation $p(w)=q(w)$, transport its left side,
reverse the resulting equation, transport its right side, and reverse
again. This proves $p(t)=q(t)$ from $s=t$ and $p(s)=q(s)$.

For a compound template $D_{x,y}E(w,x,y)$, argue by external induction
on formula depth. Open $D_{x,y}E(s,x,y)$ at a fresh $c$ using ordinary
elimination. Introduce $D_{x,y}E(t,x,y)$ at witness $c$. Since $c$ is
a fresh parameter, its syntax is not literally zero, and Intro selects
cases; no proposition $c\ne0$ is assumed. From the temporary base
$E(t,c,0)$, use the induction hypothesis and symmetry to obtain
$E(s,c,0)$. Use the opened old clause at $Sz$ to obtain $E(s,c,Sz)$,
then transport forward by the induction hypothesis to $E(t,c,Sz)$.
Close Intro and ordinary elimination. Every recursively transported
matrix has smaller formula depth. This is not induction over quantified
formulas inside the arithmetic theory.
\end{proof}

\subsection{Exact equivalence with unrestricted witness discharge}

Let $\mathcal W$ be the incoming canonical witness-discharge calculus:
the same terms, formulas, clauses, fixed axiom and Use/Intro rules, with
\[
 \frac{\Gamma\vdash D_{x,y}B\qquad
       \Gamma,\mathcal B_b\vdash b=t\qquad
       \Gamma,\mathcal B_b\vdash F}
      {\Gamma\vdash N(F[t/b])}
 \tag{Old Name}
\]
and the same freshness and export conditions. Let $\mathcal A$ denote
the deterministic calculus just specified.

\begin{theorem}[Identity-language sequent equivalence]
\label{thm:det-equivalence}
For every finite canonical clause context $\Gamma$ containing the
distinguished reflexivity axiom and every well-scoped canonical formula
$A$,
\[
               \Gamma\vdash_{\mathcal A} A
          \quad\Longleftrightarrow\quad
               \Gamma\vdash_{\mathcal W} A.
\]
The two translations are effective transformations of finite derivations
and preserve the complete end sequent literally.
\end{theorem}
\begin{proof}
From $\mathcal A$ to $\mathcal W$, translate the two premises of each
new Name recursively and supply the missing equality proof from
Lemma~\ref{lem:det-replacement}. All other rules are shared.

In the converse direction, recursively translate the three premises of
an Old Name inference. In its extended witness context, use
Corollary~\ref{cor:det-equality} to transport the translated branch
formula along the translated equation $b=t$. This gives the prescribed
normal form $N(F[t/b])$. Its original export condition says that it has
no free $b$, so Lemma~\ref{lem:det-ordinary-elim} closes the major
existential. This works even when $t$ contains $b$, because the same
ordinary equality substitution and final export condition apply.

The apparent dependency is not circular: ordinary elimination and
Cases were obtained before any equality transport; the naming-matrix
proof of term transport uses only those constructions, the supplied
premises, and one new Name. Formula transport then uses strictly
smaller matrix depth. Every replacement is therefore a finite, explicitly
defined proof construction.
\end{proof}

\begin{corollary}[Full iPRA, not just the equational fragment]
The ordinary realization of $\mathcal A$ is bi-interpretable and
equiconsistent with full first-order intuitionistic primitive recursive
arithmetic. Both numerical interpretation domains are whole, equality
is actual equality, and both numerical composite maps have graph $x=y$.
\end{corollary}
\begin{proof}
Compose Theorem~\ref{thm:det-equivalence} with the incoming canonical
calculus equivalence and its inherited scalar arithmetic and logical
bi-interpretation. The present translation is the identity on the entire
native language and every end sequent, so it preserves both inherited
formula round trips and the numerical identity maps. In particular it
preserves the designated contradiction $0=1$, yielding equiconsistency.
The scalar compiler, equality reflection, predecessor, logical macros,
and arithmetic bridge are inherited and must be included or precisely
referenced; the local result does not replace their proofs.
\end{proof}

\subsection{Implementation invariants and claim boundaries}

A forward checker may infer the major formula and the branch formula.
At Name it checks that the major has a $D$ root, opens its matrix at the
newest fresh parameter, computes its zero and first instances, and
determines $\tau_b(B)$ by~\eqref{eq:first-instance-test}. It then checks
the branch under the one new clause, performs simultaneous free-parameter
substitution, normalizes, and verifies that the result has no free $b$.
The original outer goal is compared with this inferred result.

Scope validation is essential even if repeated formula-grammar checks are
removed: a branch may use the fresh witness in an equation, but Name may
export only an outer-scoped formula. The clause antecedent must be the
actual normalized zero instance of the same opened matrix; replacing it
by a claimed true formula would be unsound. The recognized first instance
must be an equation with left side literally the current witness. Merely
having an equation root, being provably equal to that shape, or having the
witness on the right does not authorize the substitution in this rule.

There is a useful implementation invariant: if the computed zero instance
is literally $H$, then the matrix has an equation root, and so does its
first instance. Normalization never deletes a $D$ constructor. A checker
which has already validated the major formula may therefore omit a second
equation-root test on the first instance; the test $A_0=H$ already ensures
that its two equation fields exist. This is a structural fact about the
validated grammar, not a relaxation of the mathematical guard.

Numerical substitution can change the selected Intro mode or create the
first-instance guard. A primitive-tag-preserving substitution theorem is
not claimed. Admissibility follows through the exact equivalence with the
reference calculus and its established substitution theorem or repairs.

The simplification deletes one proof premise and two certificate payloads
(the replacement term and, for a synthesizing checker, the claimed branch
formula) from the earlier interface. It still has three inference families
and one explicit reflexivity clause. A measured checker reduction must
count the deterministic guard and normalization in full. It is neither a
proof of absolute minimum size nor a claim that three families first
occurred here. The new mathematical assertion is the exact completeness
of the first-instance restriction.

Uniform full normalization over arbitrary finite iterator descriptions is
total recursive but need not be primitive recursive. A finite budget or
explicit reduction certificates can give a primitive-recursive bounded
certificate predicate without changing the provable formulas; that
interface is distinct and its code must also be counted. No object-level
uniform evaluator or new totality axiom is asserted by this metasyntactic
operation.


% Self-contained normalization and raw/normal-form comparison.
% Adapted with explicit attribution from Research/full_normalization and
% incoming ipra_witness_discharge_2026-11-11/report.tex, normalization appendix.
% New explanatory addition: the administrative-let proof for the strict
% normalizer, including evaluation of an argument erased by the update.

\section{Normalization and a direct bridge to raw Choice-DI}
\label{sec:normalization-bridge}
\label{sec:canonical-bridge}

This section reconstructs the normalization and raw-comparison arguments
of the repository's full-normalization package~\cite{normalization} and the
incoming witness-discharge report~\cite{witness}. The
compound-envelope repair is inherited from those sources. The direct
comparison below connects the present deterministic Name calculus to the
raw Choice-DI foundations included later, without taking the incoming
witness-discharge equivalence as a black box. The explicit
administrative-\texttt{let} argument also accounts for the selected
normalizer's strict evaluation of a recursive value that its update may
subsequently discard.

All assertions about normalization in this section are external theorems
about finite syntax. They are not new arithmetic axioms or an
object-language uniform evaluator.

\subsection{Strong normalization of scalar computation}

On raw terms $x,0,St,I(n,a;u.h)$ let $\longrightarrow$ be the contextual
closure, including under update binders, of
\[
 I(0,a;u.h)\longrightarrow a,\qquad
 I(Sn,a;u.h)\longrightarrow h[I(n,a;u.h)/u].
 \tag{Raw computation}\label{eq:nb-reduction}
\]
Bound variables are freshened before substitution. These reductions are
stable under capture-avoiding substitution. Call a finite term strongly
normalizing if it has no infinite reduction sequence. For such a term $t$,
write $\nu(t)$ for the largest length of a reduction beginning at $t$ and
$|t|$ for its syntax-tree size. The maximum exists because the reduction
tree is finitely branching and has no infinite path.

\begin{lemma}[Closure under iteration]
\label{lem:nb-iterator}
Suppose $n,a$ are strongly normalizing and $h[v/u]$ is strongly normalizing
for every strongly normalizing $v$. Then $I(n,a;u.h)$ is strongly
normalizing.
\end{lemma}
\begin{proof}
Taking $v=u$ shows that $h$ is strongly normalizing. Induct
lexicographically on
\[
            (\nu(n),\ |n|,\ \nu(a)+\nu(h)).
\]
It suffices to show that every one-step reduct of the iterator is strongly
normalizing. A contraction inside $n$ decreases the first coordinate.
A contraction inside $a$ or $h$ decreases the last. If $h\to h'$, then
$h'[v/u]$ is a reduct of $h[v/u]$ for every $v$, so the universal body
hypothesis persists. The zero root reduct is $a$. For a successor root
$n=Sp$, we have $\nu(p)=\nu(n)$ and $|p|<|n|$. The induction hypothesis
makes $I(p,a;u.h)$ strongly normalizing, and the hypothesis about $h$
then applies to the successor root reduct
$h[I(p,a;u.h)/u]$.
\end{proof}

\begin{theorem}[Strong normalization]
\label{thm:nb-sn}
Replacing the free variables of any fixed finite raw term by strongly
normalizing terms produces a strongly normalizing term. In particular,
every finite raw term is strongly normalizing.
\end{theorem}
\begin{proof}
Induct on the fixed raw term, using simultaneous capture-free substitutions.
Variables, zero, and successor are immediate. At an iterator, first rename
its bound state variable away from the finitely many substituting terms.
The structural induction hypotheses give strongly normalizing counter and
initializer instances. The induction hypothesis for the strictly smaller
update term applies to every additional strongly normalizing assignment
for its state variable. It gives exactly the universal premise of
Lemma~\ref{lem:nb-iterator}. The identity substitution yields the last
assertion.
\end{proof}

\subsection{Confluence and normalization--substitution coherence}

\begin{lemma}[Confluence]
\label{lem:nb-confluence}
The reduction~\eqref{eq:nb-reduction} is confluent. Every term has a unique
normal form $N(t)$, up to alpha-renaming, and literally with de Bruijn
indices.
\end{lemma}
\begin{proof}
First check local confluence. Disjoint redex contractions commute, and
the zero and successor root patterns do not overlap. If an inner redex
lies inside a root contraction, contract each of its residual copies on
the other path. A zero contraction returns its initializer and erases
the entire update; erased residuals require no step. A successor
contraction may copy its recursive iterator into
several update positions. The update is finite, so only finitely many
copies arise, and substitution stability gives a common reduct after
contracting them. An inner contraction of the update itself is handled
both in the outer substituted update and in the retained recursive
iterator. These are all nested overlaps.

Strong normalization and local confluence imply confluence: induct on
the reduction height of a common source to join the first steps of two
finite reductions and then their remaining tails. Thus two normal forms
reachable from a term are equal up to bound-variable naming.
\end{proof}

A term is normal exactly when its children are normal and no iterator
counter is headed by $0$ or by $S$. Formula normalization changes only
numerical terms; the equality and $D$ constructors are preserved.

\begin{lemma}[Coherence]
\label{lem:nb-coherence}
For every finite raw term $t$ and capture-free substitution $\sigma$,
\[
 N(t[\sigma])=
 N\bigl(N(t)[x\mapsto N(\sigma x)]\bigr).
 \tag{Coherence}\label{eq:nb-coherence}
\]
The same equality holds for formulas. Normalized substitution consequently
obeys identity, composition, renaming, and the usual binder-lifting laws.
\end{lemma}
\begin{proof}
The reduction $t\to^*N(t)$ survives raw substitution. Within the finitely
many substituted copies, reduce each $\sigma x$ to $N(\sigma x)$.
The resulting terms are convertible, and confluence makes their normal
forms identical. Applying this observation twice and using composition
of raw substitutions proves
\[
 N\bigl(N(t[\sigma])[\tau]\bigr)
   =N\bigl(t[x\mapsto N((\sigma x)[\tau])]\bigr).
\]
Identity, renaming, and lifting are special cases. The formula assertion
follows structurally because formula constructors do not reduce.
\end{proof}

\paragraph{The strict implementation also terminates.}
The selected traversal normalizes the counter, initializer, and update
first. At a successor counter it normalizes the predecessor iterate
before inserting it into the update. If the update ignores its state,
that recursive value would have been erased by the direct raw root
contraction, so this strict evaluation deserves an explicit justification.

Introduce an administrative term $\mathsf{let}\ u=a\ \mathsf{in}\ h$,
with its usual capture-free reduction to $h[a/u]$, and replace the
successor root rule in the auxiliary syntax by
\[
 I(Sn,a;u.h)\longrightarrow
       \mathsf{let}\ v=I(n,a;u.h)\ \mathsf{in}\ h[v/u],
 \tag{Strict expansion}
\]
where $v$ is fresh. Reduce under all constructors and let components.
The zero rule is unchanged. There are no new normal forms: every let
has a root contraction, and its contraction composes a strict expansion
to the original successor computation.

For this auxiliary reduction, a let is strongly normalizing whenever
its value is strongly normalizing and its body sends every strongly
normalizing substitution to a strongly normalizing term. Induction on
the sum of the component reduction heights proves this: internal
contractions decrease the sum, and the root contraction is covered by
the universal body hypothesis. The proof of
Lemma~\ref{lem:nb-iterator} now uses this let-closure fact at its successor
root, after obtaining the predecessor iterate by the same lexicographic
induction. Structural induction, also including the let constructor,
therefore proves strong normalization of the auxiliary syntax.

The strict traversal is an evaluation strategy in this auxiliary
calculus: after strict expansion, it first normalizes the let value and
then contracts the let. Thus its recursive computational work cannot
continue indefinitely, even in the erased-value case. Between those
contractions it performs finite tree traversals and index renamings.
Forgetting administrative lets preserves conversion in the original
calculus; its returned normal term is therefore the unique $N(t)$.
The same argument and~\eqref{eq:nb-coherence} identify hereditary
substitution with normalized capture-free substitution. No uniform
termination theorem inside iPRA has been invoked.

\subsection{The raw reference calculus and its harmless clause convention}

Let $\mathcal C_{\mathrm{raw}}$ be raw Choice-DI over the terms of
\eqref{eq:nb-reduction} and formulas $s=t,D_{x,y}B$. It has the two raw
computation equations, formula transport, clause Use, ordinary paired
elimination, and the two separate introductions
\[
 \frac{\Gamma,B(t,0)\vdash B(t,Sz)}{\Gamma\vdash D_{x,y}B},
 \qquad
 \frac{\Gamma,E(t,z)\vdash E(t,Sz)}{\Gamma\vdash D_{x,y}E},
 \tag{Raw introductions}\label{eq:nb-raw-intro}
\]
where the second matrix is a literal equation and all eigenvariable
conditions are as in the inherited Choice-DI foundations.

For comparison we allow clauses with zero or one antecedent and retain
$\mathcal R=[\,;y.y=y]$. This is only an interface convention. In the
inherited single-antecedent raw calculus, replace an unconditional clause
$[\,;y.B]$ by $[H;y.B]$ and supply the theorem $H$ to each use.
Every use of $\mathcal R$ is replaced by derived raw reflexivity.
Conversely the optional-premise interface admits all single-antecedent
clauses. The inherited computation theorem used as a neutral antecedent
there and $H$ are both theorems, so those ordinary-assumption encodings
are interchangeable by inserting the corresponding neutral proof.
No new arithmetical or logical axiom is added by this convention.

The included Choice-DI development derives equation induction with
arbitrary passive contexts, numerical cases, all ordinary equality
reasoning, and binder congruence. Its logical and scalar arithmetic
bridges identify its ordinary realization with full first-order iPRA.
We use those results, rather than an assertion only about standard
natural-number models.

\subsection{Raw provability of normalization, including update binders}

Transport in a numerical parameter does not permit a free variable in
an equality proof to become captured by an iterator update. The needed
update congruence follows from literal-equation induction.

\begin{lemma}[Update extensionality]
\label{lem:nb-update}
If $u$ is fresh for $\Gamma$ and
$\Gamma\vdash_{\mathcal C_{\mathrm{raw}}}h(u)=k(u)$, then, after
freshening binders away from the passive data,
\[
 \Gamma\vdash_{\mathcal C_{\mathrm{raw}}}
     I(n,a;u.h)=I(n,a;u.k)
\]
for every counter $n$ and initializer $a$.
\end{lemma}
\begin{proof}
Apply derived literal-equation induction to
$E(z):=(I(z,a;u.h)=I(z,a;u.k))$, with $z$ fresh for the retained data.
At zero both sides compute to $a$. In the successor step, write the
left and right predecessor iterates as $L$ and $R$. Computation gives
$h(L)$ and $k(R)$. Capture-free numerical substitution in the supplied
pointwise proof gives $h(L)=k(L)$; the induction hypothesis $L=R$
and ordinary term transport give $k(L)=k(R)$. Transitivity and the
two successor computations prove $E(Sz)$. Induction gives $E(n)$.
Only the literal equation $E$ is inducted on.
\end{proof}

\begin{lemma}[Internal normalization equivalence]
\label{lem:nb-internal}
For every fixed finite raw term $t$,
\[
              \vdash_{\mathcal C_{\mathrm{raw}}}t=N(t).
\]
For every fixed raw formula $A$, the formulas $A$ and $N(A)$ are
interderivable, uniformly under additional retained contexts.
\end{lemma}
\begin{proof}
Each root contraction is one raw computation equation. Ordinary
transport handles a contraction inside a successor, counter, or
initializer. For a contraction beneath an update binder, use
Lemma~\ref{lem:nb-update}, repeating it for nested update contexts.
Thus every step in a finite normalization sequence is internally a
provable equality. Transitivity gives $t=N(t)$.

For formulas, induct externally on formula depth. At an equation,
normalize its two terms using the proved term equalities. At
$D_{x,y}B$, open the source formula at a fresh witness, introduce the
target at that same witness by case introduction, convert the new
zero base to the old one by the induction hypothesis, use the old
clause at a fresh successor, and convert the result back. Reverse the
conversions for the reverse implication. Substitution may create fresh
redexes in both an instance of $B$ and the corresponding instance of
$N(B)$: apply the smaller-depth induction hypothesis to both instances
and use coherence to identify their common normal form. These are uniform
derivations under every passive context. The construction uses
derived binder congruence; it never captures a free variable by an
ordinary equality transport.
\end{proof}

\paragraph{Replacing a retained clause.}
If antecedents $A,A'$ are uniformly interderivable and bodies $B(y),B'(y)$
are uniformly interderivable, then the clauses $[A;y.B]$ and
$[A';y.B']$ are interchangeable in derivations. Replace each use of the
old clause by conversion to the new antecedent, use of the new clause,
and conversion of its specialized conclusion. Induction on proof trees,
with fresh local eigenparameters, justifies this under further retained
clauses. The same statement holds for unconditional clauses, with no
antecedent conversion. In particular, a raw context and its normalized
context have the same consequences after the corresponding formula
conversions. This supplies the context-conversion step explicitly,
rather than presuming that ordinary-assumption cut alone rewrites a
conditional-universal clause.

\subsection{Repairing the omitted case introduction at witness zero}

Write $\mathcal A$ for the deterministic three-family calculus specified
in Section~\ref{sec:deterministic-name}. Its previously derived ordinary
elimination and numerical cases require no equality transport. For unused
$u$, put $q_v F(v)=D_{u,v}F(v)$. For an equation $E=(s=t)$ define
\[
 G_vE=(C(v,0,s)=C(v,0,t)),\qquad L(E)=q_vG_vE,
\]
with all terms normalized and all new variables fresh. Direct computation
gives $G_0E\equiv H$ and $G_{Sw}E\equiv E$.

\begin{lemma}[Compound envelope]
\label{lem:nb-envelope}
In $\mathcal A$, an equation $E$ and its compound envelope $L(E)$ are
interderivable under every retained context, without equality transport.
\end{lemma}
\begin{proof}
Under $E$, introduce the unused witness of $q_vG_vE$ at $1$. Its
successor premise is literally $E$, so this is case introduction.
Conversely, open $q_vG_vE$ by derived ordinary elimination. Its clause
has antecedent $H$. Prove $H$ by the fixed reflexivity clause and use
the opened clause at $1$ to obtain $E$.
\end{proof}

\begin{lemma}[Unrestricted case introduction]
\label{lem:nb-all-cases}
The case introduction in~\eqref{eq:nb-raw-intro}, with normalized
substitution throughout, is derivable in $\mathcal A$ for every witness
and every matrix.
\end{lemma}
\begin{proof}
Only a literal-equation matrix $E(x,y)$ at the literal witness zero
is missing from the primitive case mode. Suppose
$\Gamma,E(0,0)\vdash E(0,Sz)$ is supplied. Put
$A(x,y)=L(E(x,y))$ with suitably fresh envelope binders. Under
$A(0,0)$, unwrap its envelope, insert the supplied derivation by
ordinary-assumption cut, and wrap the result. This gives $A(0,Sz)$.
Since $A$ has a $D$ root, primitive Intro at zero now selects cases,
and yields $D_{x,y}A(x,y)$.

Open it by ordinary elimination at a fresh $c$. Introduce
$D_{x,y}E(x,y)$ at witness $c$. Its syntax is a parameter rather than
literal zero, so Intro again selects cases. Under $E(c,0)$, wrap the
equation to $A(c,0)$, use the opened clause at $Sz$, and unwrap
$A(c,Sz)$ to $E(c,Sz)$. Closing Intro and ordinary elimination gives
the desired formula. No claim that $c\ne0$ is used.
\end{proof}

\begin{lemma}[Equation progress at every witness]
\label{lem:nb-all-progress}
The equation-progress introduction in~\eqref{eq:nb-raw-intro}, with
normalized substitution, is derivable in $\mathcal A$ for every witness.
\end{lemma}
\begin{proof}
Given $\Gamma,E(t,z)\vdash E(t,Sz)$, use a fresh unused witness
coordinate $w$. Primitive Intro at zero on the equation matrix
$E(t,y)$ yields $D_{w,y}E(t,y)$ by progress. Under the prospective
outer case assumption $E(t,0)$, open this formula by ordinary
elimination and use its clause at $Sz$. This provides the successor
premise for unrestricted case introduction of $D_{x,y}E(x,y)$
at witness $t$, available by Lemma~\ref{lem:nb-all-cases}.
The unused coordinate prevents the opened witness from replacing $t$.
\end{proof}

\subsection{An auxiliary normal-form calculus}

Let $\mathcal C_{\mathrm{nf}}^+$ use canonical terms and formulas,
normalized substitution, optional-premise clauses and the fixed
reflexivity clause. Its rules are Use, ordinary paired elimination,
full formula transport, and the separate unrestricted case and
literal-equation progress introductions. There is no separate
computation axiom, since both computation equations normalize to
reflexivity.

\begin{proposition}[Canonical Choice-DI and the new calculus]
\label{prop:nb-canonical-choice}
The calculi $\mathcal C_{\mathrm{nf}}^+$ and $\mathcal A$ derive the
same canonical sequents.
\end{proposition}
\begin{proof}
Translate $\mathcal C_{\mathrm{nf}}^+$ to $\mathcal A$ by using its
derived ordinary elimination, term and formula transport, and
Lemmas~\ref{lem:nb-all-cases}--\ref{lem:nb-all-progress}. Use and the
fixed reflexivity clause are shared.

Conversely each primitive Intro is one of the two available choices.
For Name, ordinary elimination opens its witness clause. When the
deterministic guard succeeds, Use at $1$, supplied with reflexivity
of zero, proves the selected equality; otherwise reflexivity proves
$b=b$. Full formula transport gives the exported normalized formula,
whose stipulated freedom from $b$ permits ordinary elimination to
close. All translations preserve the entire canonical end sequent.
\end{proof}

\subsection{Both proof translations and both composites}

Let $T$ normalize all numerical terms in raw formulas and clauses.
Let $U$ regard canonical syntax as raw syntax. These are operations on
finite descriptions, not extra operations on the numerical domain.

\begin{theorem}[Direct raw/canonical comparison]
\label{thm:direct-raw-canonical}
For every raw sequent and every canonical sequent,
\begin{align*}
 \Gamma\vdash_{\mathcal C_{\mathrm{raw}}}A
  &\quad\Longleftrightarrow\quad
        T\Gamma\vdash_{\mathcal A}T A,\\
 \Delta\vdash_{\mathcal A}B
  &\quad\Longleftrightarrow\quad
        U\Delta\vdash_{\mathcal C_{\mathrm{raw}}}U B.
\end{align*}
Moreover $TU(B)=B$ literally, and $A$ and $UT(A)$ are uniformly
interderivable in the raw calculus. The corresponding clause-context
composites are interchangeable in every derivation.
\end{theorem}
\begin{proof}
First translate a raw derivation into $\mathcal C_{\mathrm{nf}}^+$
by induction on derivations. The two computation equations have
identical normal forms on their two sides, so become reflexivity.
Coherence~\eqref{eq:nb-coherence} makes the other normalized
premises and conclusions instances of the corresponding canonical
rules. Equation-root status is preserved, and normalization cannot
introduce a new free eigenparameter. Crucially, the auxiliary calculus
retains both introduction choices, so this translation never assumes
that a nonzero-syntax raw witness stays nonzero after normalization.
Proposition~\ref{prop:nb-canonical-choice} then translates to
$\mathcal A$.

In the reverse direction first translate $\mathcal A$ to
$\mathcal C_{\mathrm{nf}}^+$. Regard its canonical formulas as raw.
For a raw rule whose substitution instance is not already normal,
Lemma~\ref{lem:nb-internal} converts the normalized premises to raw
instances, the raw rule is applied, and its conclusion is converted
back. The retained clause of an elimination may have a normalized
zero antecedent instead of the literal raw substitution instance;
the explicit clause-replacement construction above converts that
context. The same reasoning handles Intro's temporary ordinary
assumption. Raw reflexivity supplies the fixed reflexivity-clause
uses. Thus induction on canonical derivations gives a raw proof
with the stated end sequent.

On canonical syntax $TU$ is literal identity. On raw syntax $UT$
is formula normalization, whose internal equivalence and clause
interchangeability were proved above. Composing the two proof
translations with these equivalences gives both reverse implications
in the displayed statements as well as the composites.
\end{proof}

\begin{corollary}[Full first-order bi-interpretation]
\label{cor:nb-full-ipra}
Composing the direct comparison with the included Choice-DI logical and
scalar arithmetic bridges gives a bi-interpretation of the present
calculus with full first-order intuitionistic primitive recursive
arithmetic. Both numerical domains are whole, both equalities are
actual equality, and both numerical composite isomorphisms have graph
$x=y$. Both complete formula composites are internally provable.
The theories are equiconsistent.
\end{corollary}
\begin{proof}
The inherited bridge reconstructs all intuitionistic logical
connectives and both quantifiers, quantifier-free induction, and the
fixed-description primitive-recursive function symbols, with both
formula round trips. Theorem~\ref{thm:direct-raw-canonical} supplies
the remaining change of term presentation. Its operations commute
with substitution modulo internally provable computation equalities,
so it preserves those interpretations. Numerical variables normalize
to themselves: the comparison changes finite term descriptions,
not the numerical domain or equality. Hence the inherited numerical
identity graphs remain $x=y$.

The common designated contradiction $0=S0$ is already normal and
is preserved in both directions. The full formula comparisons
therefore imply equiconsistency, as well as the stronger
bi-interpretability statement.
\end{proof}

\subsection{Uniform normalization is external}

The full first-order iPRA convention here has a symbol for each fixed
primitive-recursive definition and the corresponding defining
equations, with quantifier-free induction in intuitionistic logic.
It is neither the equation-only fragment nor a theory with unrestricted
formula induction. A fixed scalar term compiles to one such
primitive-recursive definition; no numerical variable ranges over all
definitions inside the object theory.

Uniform normalization over arbitrary finite scalar-term codes is total
recursive but is not primitive recursive. Otherwise, by coding a unary
scalar term, substituting a numeral, normalizing, and extracting the
resulting numeral, one would obtain a primitive-recursive universal
evaluator $V(e,n)$ for all primitive-recursive functions. Invalid codes
could be assigned zero. The diagonal function $V(e,e)+1$ would itself
have a scalar-term code, contradicting universality. This uses the
included scalar compilation theorem, not an assumed uniform evaluator.

The size of a complete checker must count its normalizer, and the
unbounded interface has no primitive-recursive running-time bound.
Explicit finite fuel or reduction traces give a different,
primitive-recursive bounded certificate interface for the same
theorems. They do not add an arithmetic assertion of a uniform bound
on the necessary fuel. The metatheorems here are written proofs; the
finite executable checks are implementation evidence rather than a
new Lean or Rocq verification of them.


\section{Complete checker, certificate syntax, and measurements}
\label{sec:implementation}

The trusted executable is one Python module, \path{code/checker.py}.
It imports only \texttt{json} and \texttt{sys} from the standard library.
Proof construction, translation, measurement, and testing are outside the
acceptance kernel. The complete normalized module is printed in
Appendix~\ref{app:checker}. Standard Python and its JSON implementation
are the common host environment, as in the repository's comparisons.

\subsection{Numerical and formula codes}

The integer $0$ is numerical zero. Negative integers denote free
parameter levels; positive integers denote bound de Bruijn indices,
starting with $1$ for the nearest binder. Successor is the two-entry
array \texttt{[0,t]}, and iteration is \texttt{[n,a,h]}, binding one
variable in \texttt{h} only. An equation is \texttt{[s,t]}; a paired
formula is \texttt{[B]}, binding two variables in \texttt{B}. Sorting
distinguishes the successor array from an equation. The walker checks
the original syntax before returning its normalized form, so malformed
material in an erased iterator body is not accepted.

At free scope $m$ and binder depth $d$, an integer numerical leaf must
satisfy $-m\le i\le d$. Boolean values are not integers for this test.
An iterator adds one to $d$ in its body; a $D$ adds two. Opening the
outer witness of $D$ removes its index $2$ while preserving the inner
control index $1$. Bound substitution shifts an inserted term under
internal binders. In contrast, a normalized target obtained from the
fully opened first clause instance has no loose positive indices;
free-parameter replacement may insert it without an additional shift.

Contexts are Python lists of pairs \texttt{(base,body)}. An empty
array \texttt{base=[]} means no antecedent. The initial context is
\texttt{[([],[1,1])]}, the exact code of the fixed reflexivity clause.
Each Intro or Name extends the context by one clause and introduces one
numerical parameter. Thus a context of length $n$ has free scope $n-1$,
and its next parameter is $-n$. No additional scope counter is stored.

\subsection{A uniform three-field proof packet}

The checker synthesizes the proved formula. Every used proof packet has
three fields:
\[
                   [\mathrm{code},\mathrm{arg},\mathrm{branch}].
\]
The three alternatives are disjoint.

\begin{center}
\begin{tabular}{@{}lll@{}}
\toprule
Packet & Meaning & Supplied data \\
\midrule
\texttt{[j,t,p]} & Use & clause index, target, antecedent proof \\
\texttt{[B,t,p]} & Intro & nonempty matrix, witness, branch proof \\
\texttt{[[],p,q]} & Name & major proof, branch proof \\
\bottomrule
\end{tabular}
\end{center}

In Use, $j$ is an exact nonnegative integer below the context length.
If the selected clause has no antecedent, its third field is unused;
\texttt{null} is the suggested value. This is an explicit certificate
convention, shared with the control. Intro's nonempty first field is a
matrix well scoped under the two binders of its resulting $D$ formula.
The empty array is not a formula and selects Name. The latter infers its
matrix from its major proof and its local formula $F$ from its branch.
It has no separate matrix, formula, replacement-term, or equality-proof
payload.

The command-line input has exactly two keys, \texttt{goal} and
\texttt{proof}. The goal must be a closed canonical formula and must
equal the inferred conclusion literally. Successful checking prints
\texttt{ACCEPT} and exits zero; every caught failure prints
\texttt{REJECT} and exits nonzero. The trusted implementation contains
no \texttt{assert} statements, so \texttt{python -O} does not disable
checks. The optional fuel interface is described separately below.

\subsection{The accepted-trace invariant}

For a well-formed context of length $n$, every successful call to
\texttt{derive} returns a canonical formula scoped over parameters
$-1,\ldots,-(n-1)$ and a valid derivation of that formula.
The proof is induction on the successful finite call tree.

Use validates its target, checks its antecedent when present, and returns
the normalized instance of the selected valid clause. Intro checks the
successor branch under its constructed ordinary assumption and returns
the paired matrix. The final syntax check validates that original matrix
including material which normalization might otherwise erase. This
late check is sufficient for the accepted-trace induction: once a call
has succeeded, the branch's context is known to have been well formed.

At Name, unpacking the major's single child requires a $D$ root. The
zero instance supplies the antecedent of precisely the clause being
opened. If that antecedent is literally $[0,0]$, the matrix has equation
root; numerical substitution and normalization preserve formula roots,
so its first instance is also an equation. A separate length test for
that row would therefore be redundant. The mapper rewrites the current
fresh leaf only when it is also the first row's left side and the base
is $[0,0]$. This is exactly the deterministic test in
\eqref{eq:first-instance-test}. The final scope check rejects a witness
which survives in the exported normal form. Inserted targets are not
visited again, so the substitution is simultaneous.

The converse direction follows by encoding the three inference families
recursively. Physical Python resource limits can reject a sufficiently
large finite derivation; as usual, mathematical completeness concerns
the specified finite-tree algorithm with adequate resources.

\subsection{Measured comparisons and their interpretation}

All reported normalized measurements use CPython 3.12.14: remove
documentation strings, apply \texttt{ast.unparse}, and append one final
line feed. Significant tokens exclude comments, indentation tokens,
and line/end markers. AST nodes are counted by \texttt{ast.walk}.
Gzip uses maximum level $9$ and timestamp zero and is informational.
Every acceptance helper, parser entry, scope and syntax check,
normalization step, and final acceptance test in the module is counted.

\begin{center}
\small
\begin{tabular}{@{}lrrrr@{}}
\toprule
Complete module & Bytes & Tokens & AST & Gzip-9 \\
\midrule
New deterministic Name, packets & 2157 & 704 & 804 & 870 \\
Matched unrestricted Name, packets & 2228 & 701 & 813 & 894 \\
Exact incoming witness discharge & 2273 & 732 & 813 & 903 \\
Exact incoming uniform binding & 2238 & 693 & 764 & 866 \\
New deterministic Name, with fuel & 2340 & 752 & 863 & 936 \\
Matched unrestricted Name, with fuel & 2411 & 749 & 872 & 959 \\
\midrule
Existing native public-goal profile & 788 & 358 & 413 & 460 \\
Existing compact canonical profile & 516 & 265 & 306 & 332 \\
\bottomrule
\end{tabular}
\end{center}

The matched packet control uses the same syntax, normalizer, validator,
context representation, source conventions, and input/output contract.
Only Name differs: it receives an equality branch as well as a formula
branch. The new restriction saves $71$ normalized bytes ($3.19\%$),
$9$ AST nodes, and $24$ gzip bytes; it costs $3$ significant tokens.
This is an actual checker-size reduction with an explicit coordinate
tradeoff, not a claim of coordinatewise dominance over the matched control.

Against the exact incoming witness-discharge module, the new combination
saves $116$ normalized bytes ($5.10\%$), $28$ tokens, $9$ AST nodes and
$33$ gzip bytes. That comparison includes the certificate-interface
change and should not be attributed entirely to the mathematical premise
deletion. Against the incoming uniform-binding module, it saves $81$
bytes but has $11$ more tokens and $40$ more AST nodes. The two compact
historical profiles remain substantially smaller; their source sizes
are not new records of this package. No claim that this construction
beats the repository's global lower envelope is justified by these data.

The measured raw readable source of the selected checker is $2544$
bytes, or $1059$ bytes with gzip-9. Complete source hashes and every
normalized and compressed module are supplied with the measurements.

\subsection{Executed evidence}

The executable replay suite passed $16$ positive cases and $57$ adversarial
rejection cases, with $178$ standalone process checks across ordinary
and optimized Python. Thirteen positive cases start from inherited
proofs, pass through the separate-rule control, and are translated into
the new calculus. All sixteen selected proofs are decoded, checked by an
independent annotated reference implementation, and translated back to
the old witness-discharge checker with exactly the same end formula.

The positive cases include symmetry from a retained assumption,
transitivity using both equality premises, repeated and nested term
occurrences, substitution below formula binders, equation induction that
uses its induction hypothesis, and the zero-witness case repair that
uses its base. New boundary cases exercise default discharge, the
literal-$H$ restriction, and an active self-containing replacement whose
fresh occurrences disappear after normalization. Rejection cases cover
the wrong equation source, an unsuitable base, escaped witnesses,
malformed erased syntax, ill-scoped terms, wrong JSON types, and malformed
proof packets. The machine-readable log identifies individual cases.

These checks provide executable evidence for the implementation and
translations. The universal soundness, completeness, normalization, and
bi-interpretation claims are the written mathematical theorems in this
report. This delivery does not claim a new Lean or Rocq verification.


% Mathematical note for inclusion in the parent report.
\section{A primitive-recursive certificate interface}
\label{sec:fuel-interface}

Full normalization is a terminating operation on well-formed scalar-iterator
terms, but its unbounded running time should not be confused with a
primitive-recursive bound uniform in arbitrary input descriptions. The
optional bounded checker makes the certificate relation primitive recursive
without requiring such a normalization bound.

\paragraph{Exact interface.}
The input object has exactly three fields, \texttt{goal}, \texttt{proof}, and
\texttt{fuel}. The first two have the same meanings as in the unbounded
checker. The last is an exact nonnegative integer; in particular, Boolean
values are rejected rather than treated as the integers zero and one. The
budget is reset at the start of each call to \texttt{check}. Every entry to
\texttt{walk} or \texttt{derive} first requires a positive remaining budget
and then subtracts one. Exhaustion rejects the certificate. No other
mathematical rule, term operation, scope check, or normalization convention
changes. The same transformation is applied to the matched old-discharge
control.

The budget counts recursive entries, not seconds, bytes, reductions, or
Python instructions. A successful run may leave unused fuel. This explicit
interpretation is necessary when comparing the new and control checkers:
one may perform more normalization traversals even when its proof interface
contains fewer premises.

\paragraph{Soundness and preservation of the theorem set.}
Erase the calls that decrement and test the budget. Every accepting bounded
run then becomes an accepting run of the corresponding unbounded checker
on the same goal and proof. Thus adding fuel cannot introduce a theorem.
Conversely, consider any particular successful finite run of the unbounded
checker. Let $k$ be the number of entries to its two metered procedures. A
budget of exactly $k$ follows that same run and accepts; every larger budget
does so as well. A budget of $k-1$ rejects at the last required entry. Hence
existentially quantifying the fuel field preserves the accepted theorems:
\[
  \exists p\;\operatorname{Accept}(A,p)
  \quad\Longleftrightarrow\quad
  \exists p\,\exists k\;\operatorname{Accept}_{\mathrm{fuel}}(A,p,k).
\]
This argument concerns each fixed finite derivation. It does not assert a
primitive-recursive bound on the fuel needed by all proofs, nor does it assert
that normalization itself has such a uniform bound.

\paragraph{Why bounded acceptance is primitive recursive.}
Use an ordinary finite-tree encoding of parsed input data, stack frames,
closures from the fixed program, and finite lists. Parsing finite JSON and
the initial exact-key and integer tests can be implemented by bounded
search and primitive recursion on the input string.

Every cycle in the checker call graph passes through \texttt{walk} or
\texttt{derive}. The other helpers either perform a finite local operation
or immediately enter one of these procedures. The callbacks supplied to
\texttt{walk} come from a fixed finite collection of program expressions;
the certificate cannot supply executable callbacks. In particular, a shift
requested by a callback enters \texttt{walk} and consumes fuel. A run with
initial fuel $k$ therefore has at most $k$ successful metered activations,
with an additional failing entry possible on exhaustion. Unwinding those
activations cannot restart an unmetered recursive computation.

The remaining local operations are integer arithmetic and comparison,
finite list construction and access, list or tree equality, and traversal
of a finite input collection. Each has a primitive-recursive implementation.
Tree sizes do require an explicit bound, because substitution can duplicate
an existing subtree. Such duplication does not invalidate primitive
recursiveness: a single local constructor or substitution-return operation
has output length bounded by a fixed polynomial in the total length of the
already available finite data. List comparisons inspect only those data.
The number of stack-frame entries and returns is bounded by a constant
multiple of $k+1$, and all completed iterations of an explicit traversal
either inspect input data or enter a metered child.

More conservatively, let $N$ bound the length of the input encoding and let
$b_0=N+k+10$. Iterating an elementary majorant, for example
\[
             b_{j+1}=2^{(b_j+1)^2},
\]
a fixed constant multiple of $k+N+1$ times bounds the finite data and
continuations generated by this fixed program. The constant is chosen
larger than its fixed number of local operations per activation; it does
not depend on a certificate. A further elementary bound covers the local
list comparisons and integer operations. The resulting total bound is
primitive recursive in $(N,k)$, by primitive recursion on the displayed
iteration count. One may consequently simulate the defunctionalized
checker for that many elementary steps and return its acceptance bit.
This establishes a primitive-recursive predicate for bounded acceptance.

The bound is deliberately very loose. No such majorant is evaluated by the
delivered Python checker: it simply maintains the integer budget. The
argument explains why its abstract bounded execution has a
primitive-recursive arithmetization, rather than assuming that bounded
normalization is free.

\paragraph{Implementation boundary.}
As with the unbounded reference implementation, a concrete Python process
can reject a certificate because of a finite stack, memory, parser, or
execution limit. The finite-derivation completeness assertion above uses
the abstract algorithm with adequate resources. Such implementation limits
can only add rejection; they do not supply an additional inference rule.
The fuel field is an explicit mathematical bound and is independent of
those practical limits.

\paragraph{Executed boundary checks.}
The accompanying validation script tests five accepted certificates for
both the new and matched control checkers: closed reflexivity, a default
discharge with the wrong row source, an active self-containing target whose
occurrences disappear after normalization, induction that uses its
hypothesis, and substitution beneath a formula binder. For each checker
and certificate it measures the exact metered cost, verifies acceptance at
that cost, verifies rejection at one less and at zero, rejects eight
malformed budget values, and checks the exact object-key contract. The
exact-budget and one-less outcomes are also tested through the standalone
command-line entry point. The machine-readable results and the full-module
size measurements are in \texttt{fuel\_verification.json}.


\clearpage
\appendix
\section*{Attribution and role of the inherited foundations}
\addcontentsline{toc}{section}{Attribution and role of the inherited foundations}

The following five sections reproduce the mathematical reference fragments
from the pinned \texttt{certified\_zero} package~\cite{baseline}.
They retain their original attributions and theorem dependencies.
Words such as ``new'' and ``old'' within those inherited fragments refer
to their original comparison, not to the present first-instance rule.
The raw clause convention in the first fragment represents an ordinary
assumption through a fixed true antecedent; an unconditional constant
clause gives exactly the same assumption, since that antecedent is a
theorem. The canonical comparison in
Section~\ref{sec:canonical-bridge} supplies the passage to the present
syntax and initial-clause convention.

Their purpose is substantive: they derive the ordinary logical rules,
coverage, predecessor, reflection of actual equality, quantifier-free
induction, pairing, and primitive recursion, and prove both numerical
and formula composites internally. The new local theorem depends on
that bridge and does not claim to replace its arithmetic construction.


% New mathematical result: complementary introduction choices.
% Requires amsmath/amssymb/amsthm and theorem, lemma, proposition environments.
% The inherited arithmetic-to-iPRA bridge must be supplied separately.

\section{Complementary introductions with one local assumption}
\label{sec:choice-introduction}

The old paired introduction retains both the zero instance and an
induction instance when its matrix is an equation. Both together are
unnecessary. A certificate may choose either one, provided the induction
choice is restricted to a literal equation. This changes neither the
formula language nor the meaning of a formula. The omitted combination
is recovered by a short nested introduction--elimination construction.

\subsection{Complete native syntax and rule inventory}

The scalar numerical language is
\[
 t::=x\mid0\mid St\mid I(n,a;u.h).
\]
Only $h$ binds the state variable $u$; $n$ and $a$ do not.
The update may contain fixed numerical parameters. The computation
axioms in every retained context are
\[
 I(0,a;u.h)=a,\qquad
 I(Sn,a;u.h)=h[I(n,a;u.h)/u].
 \tag{I0, IS}\label{eq:choice-computations}
\]
The argument also applies unchanged to the full fixed-description PR
language with its projection, composition, and recursion equations.
No numerical variable ranges over function descriptions, and no uniform
evaluator is postulated.

Native formulas and reusable context clauses are
\[
 A::=(s=t)\mid\mathsf D_{x,y}A,\qquad
 \langle A;y.B(y)\rangle,\quad y\notin\operatorname{FV}(A).
\]
Bound variables are renamed before capture-free substitution.
The two variables $x,y$ in a paired formula are distinct.
An ordinary assumption $F$ is the constant clause
$\langle H;v.F\rangle$, where $v$ is absent from $F$ and $H$ is the
fixed computation theorem $I(0,0;u.0)=0$. Thus ordinary hypothesis
use is a clause Use preceded by this computation proof.

There are computation axioms and exactly the following inference
patterns:
\[
\frac{\Gamma\vdash A}{\Gamma\vdash B(t)}
\quad\mathsf{Use}
\quad(\langle A;y.B(y)\rangle\in\Gamma),
\]
\[
\frac{\Gamma\vdash\mathsf D_{x,y}B(x,y)\qquad
      \Gamma,\langle B(a,0);y.B(a,y)\rangle\vdash F}
     {\Gamma\vdash F}
\quad\mathsf{DE},
\qquad
\frac{\Gamma\vdash s=t\qquad\Gamma\vdash F(s)}
     {\Gamma\vdash F(t)}
\quad\mathsf{Tr},
\]
\[
\frac{\Gamma,B(t,\eta_\varepsilon(z))\vdash B(t,Sz)}
     {\Gamma\vdash\mathsf D_{x,y}B(x,y)}
\quad\mathsf{DI}_{\varepsilon},
\qquad
\eta_0(z)=0,\quad\eta_1(z)=z.
\tag{Choice DI}\label{eq:choice-di}
\]
Here $\varepsilon$ is a certificate choice in $\{0,1\}$, and
$\varepsilon=1$ is permitted only when $B$ has a literal equation
root. An equation beneath a binder is not a literal equation.
Call $\varepsilon=0$ the \emph{case choice} and $\varepsilon=1$
the \emph{progress choice}. Both retain exactly one ordinary
formula assumption. Compound matrices have only the case choice.

In DI the eigenparameter $z$ is fresh for $\Gamma$, the witness
$t$, and every fixed parameter of $B$. In DE the eigenparameter
$a$ is fresh for $\Gamma$, the conclusion $F$, and every fixed
parameter of $B$. In transport, $F(v)$ is any fixed native
formula template. Displayed local assumptions and clauses are
discharged. There are no primitive reflexivity, coverage,
separation, ordinary logical, or additional induction rules.

We call this calculus $\mathcal C$. The original comparison
calculus $\mathcal N$ has identical syntax, computation,
Use, DE, and transport, but its introduction is
\[
\frac{\Gamma,B(t,0),[B(t,z)]_{\mathrm{at}}\vdash B(t,Sz)}
     {\Gamma\vdash\mathsf D_{x,y}B(x,y)}.
\tag{Old DI}\label{eq:choice-old-di}
\]
The bracket supplies $B(t,z)$ exactly at a literal equation.

\subsection{Exact equivalence and the unused-witness construction}

\begin{lemma}[Structural admissibility]
\label{lem:choice-structural}
Weakening, exchange, contraction, renaming, ordinary-assumption cut,
and capture-free numerical substitution are admissible in $\mathcal C$.
\end{lemma}
\begin{proof}
Induct on derivations. Rename local eigenparameters before inserting
proofs or terms. Numerical substitution preserves a formula's root
constructor, so an allowed progress choice remains allowed.
The selected choice does not change under substitution or cut.
For ordinary-assumption cut, replace each use of its constant clause
by the supplied proof. Computation schemas are closed under
capture-free numerical substitution. These facts verify every case.
\end{proof}

\begin{theorem}[Identity-language equivalence]
\label{thm:choice-equivalence}
For every retained clause context $\Gamma$ and native formula $F$,
\[
 \Gamma\vdash_{\mathcal C}F
 \quad\Longleftrightarrow\quad
 \Gamma\vdash_{\mathcal N}F.
\]
The transformations are effective on finite derivations and leave
all numerical and formula syntax unchanged.
\end{theorem}
\begin{proof}
Each new introduction is an instance of the old one after weakening
by the omitted local assumption. For the case choice at an equation,
add the induction instance. For the progress choice, add the zero
instance. Compound case introduction is already the old rule.
Thus every new derivation translates to an old one.

For the reverse direction, transform an old derivation recursively.
Only its atomic DI case needs a construction. Suppose its matrix is
the equation $E(x,y)$ and its transformed premise is
\[
 \pi:\quad\Gamma,E(t,0),E(t,z)\vdash E(t,Sz).
\]
The outer new inference will be case DI with witness $t$. It is
therefore enough to prove $E(t,Sz)$ under
$\Gamma,E(t,0)$.

Rename the eigenvariable in $\pi$ to a fresh $v$. Choose a new
variable $u$ absent from $E(t,y)$ and all retained data. Under
$\Gamma,E(t,0)$, progress DI with witness zero introduces
\[
             \mathsf D_{u,y}E(t,y)
\]
from exactly the renamed premise $\pi[v/z]$.
The bound coordinate $u$ is deliberately unused.

Apply DE to this paired formula. Because $u$ is unused, the opened
clause is exactly
\[
               \langle E(t,0);y.E(t,y)\rangle;
\]
it contains no new existential witness. The retained ordinary
assumption $E(t,0)$ supplies its antecedent. Use at $Sz$ gives
$E(t,Sz)$, and DE closes. Finally the outer case DI introduces
the original $\mathsf D_{x,y}E(x,y)$ with witness $t$.
The original premise proof $\pi$ has been inserted once.

An old compound introduction already is the new case choice.
All other rule cases are unchanged. This completes the induction
on the original derivation.
\end{proof}

The unused inner coordinate is essential. Introducing
$\mathsf D_{x,y}E(x,y)$ internally and then eliminating it would
give an arbitrary fresh witness, not necessarily the original
$t$. The construction instead keeps $t$ as a fixed parameter
of $E(t,y)$ and binds only the unused coordinate $u$.

\begin{proposition}[Linear bound on inference nodes]
\label{prop:choice-node-bound}
Count each computation axiom, Use, DI, DE, and transport as one
inference node, expanding ordinary-assumption use as a computation
of $H$ followed by Use. If an old derivation has $N$ nodes and
$r$ atomic DI nodes, the displayed transformation has exactly
\[
                   N+5r\leq 6N
\]
nodes.
\end{proposition}
\begin{proof}
Each atomic old DI node is replaced by six nodes: outer case DI,
inner progress DI, DE, Use of the opened clause, Use of the
retained base assumption, and the computation of $H$ needed for
that last ordinary-assumption use. Its old premise subtree
occurs once. All other inference nodes are copied once.
\end{proof}

This counts inference-tree nodes. It is not a bound on serialized
JSON bytes, formula sizes, checker runtime, or the number of
consumed packet cells when hypothesis closure is implicit.
In the formula-context Q/D presentation, the corresponding
replacement has outer case DI, inner progress DI, DE, QE,
and two Hyp nodes, giving the same six-node count.

\subsection{Induction and numerical cases separate cleanly}

\begin{lemma}[Contextual equation induction]
\label{lem:choice-ei}
The progress choice, DE, and Use derive
\[
\frac{\Gamma\vdash E(0)\qquad
      \Gamma,E(z)\vdash E(Sz)}
     {\Gamma\vdash E(t)}
\tag{EI}\label{eq:choice-ei}
\]
for every equation $E$, with the usual freshness condition.
\end{lemma}
\begin{proof}
Choose an unused witness coordinate $u$. Progress DI introduces
$\mathsf D_{u,y}E(y)$ from the given step. DE opens the clause
$\langle E(0);y.E(y)\rangle$. Supply its antecedent using the given
base proof and use the clause at $t$. Close DE.
\end{proof}

\begin{lemma}[Cases, equality, and binder congruence]
\label{lem:choice-preinduction}
The case choice derives, for arbitrary native formulas,
\[
\frac{\Gamma\vdash F(0)\qquad\Gamma\vdash F(Sz)}
     {\Gamma\vdash F(t)}
\tag{NC}\label{eq:choice-nc}
\]
with the usual freshness condition. Equality reflexivity,
symmetry, transitivity, numerical congruence, and uniform
matrix interderivability under $\mathsf D$ are derivable
without using the progress choice.
\end{lemma}
\begin{proof}
For cases, use case DI with an unused witness coordinate and
matrix $F(y)$, ignoring its retained base in the successor
premise. DE opens $\langle F(0);y.F(y)\rangle$; the given zero
proof and Use at $t$ finish.

For reflexivity put $r=I(0,t;u.0)$. Use the computation equation
$r=t$ as both transport premises with motive $v=t$ to obtain
$t=t$. The usual explicit equality motives give symmetry,
transitivity, and numerical congruence.

For binder congruence, open the old paired formula at a fresh
witness $a$ and introduce the new one by case DI at $a$.
Convert its retained base to the old base, use the opened
old clause at the successor, and convert back. Ordinary cut
closes the construction, and the reverse conversion gives
the converse. No induction is used.
\end{proof}

The primitive progress choice is still restricted to equations.
Allowing it at arbitrary compound matrices would, by the
argument for \eqref{eq:choice-ei}, supply induction for
arbitrary formulas and change the intended theory.

\subsection{The entire ordinary logical interface}
\label{subsec:choice-logic}

Here is a predecessor-free reconstruction of the ordinary logical
operations. It makes explicit that the case choice supplies the
logical interface independently of progress induction.

Let $C(v,a,b):=I(v,a;u.b)$ with $u$ unused. Computation gives
$C(0,a,b)=a$ and $C(Sv,a,b)=b$. Define the gate on fixed finite
formula syntax by
\[
 G_v(s=t):=(C(v,0,s)=C(v,0,t)),\qquad
 G_v(\mathsf D_{x,y}B):=\mathsf D_{x,y}G_vB,
\]
renaming bound variables away from the control. Equality,
binder congruence, and external structural induction give
\[
 \vdash G_0A,\qquad G_{Sv}A\dashv\vdash A,\qquad
 A\dashv\vdash\mathsf D_{u,y}G_yA
 \quad(u,y\notin\operatorname{FV}(A)).
\tag{Gate and padding}\label{eq:choice-gate}
\]
For padding, case DI proves the forward direction. DE, the
zero-gate theorem, and Use at $1$ prove the reverse direction.

Define a selector between two arbitrary formulas by
\[
 M_v(s=t,u=w):=(C(v,s,u)=C(v,t,w))
\]
and, with aligned fresh paired variables, by
\[
\begin{aligned}
 M_v(\mathsf D_{x,y}A,\mathsf D_{x,y}B)
   &:=\mathsf D_{x,y}M_v(A,B),\\
 M_v(E,\mathsf D_{x,y}B)
   &:=\mathsf D_{x,y}M_v(G_yE,B),\\
 M_v(\mathsf D_{x,y}A,E)
   &:=\mathsf D_{x,y}M_v(A,G_yE),
\end{aligned}
\]
where $E$ is an equation. Every call decreases maximum input
binder depth. Padding and congruence prove
$M_0(A,B)\dashv\vdash A$ and
$M_{Sv}(A,B)\dashv\vdash B$.

Put $\delta(v)=C(v,1,0)$ and $p(z)=I(z,0;u.\delta(u))$.
Only the fixed computations $p(1)=1$ and $p(2)=0$ are needed.
The full logical macro inventory is
\[
\begin{aligned}
 A\Rightarrow_*B&:=\mathsf D_{u,z}M_z(A,B),\\
 A\wedge_*B&:=\mathsf D_{u,z}G_zM_{p(z)}(A,B),\\
 \exists_*x\,A(x)&:=\mathsf D_{x,y}G_yA(x),\\
 A\vee_*B&:=\exists_*x\,M_x(A,B),\\
 \forall_*x\,A(x)&:=
        A(0)\wedge_*\mathsf D_{u,z}G_zA(z),\\
 \bot_*&:=(0=1).
\end{aligned}
\tag{Logical macros}\label{eq:choice-macros}
\]
Newly bound variables are fresh; $u$ is unused.

Implication introduction obtains its assumption from the zero
selector and supplies the consequent at successors, then uses
case DI. Elimination opens its clause and uses it at $1$.
Existential introduction uses its witness and the successor
gate; elimination opens at a fresh witness, supplies the
zero-gate theorem, and uses the clause at $1$.

Disjunction introduction uses witnesses $0$ and $1$.
For elimination, discharge the two branch proofs as
$A\Rightarrow_*F$ and $B\Rightarrow_*F$. NC applied to
$M_z(A,B)\Rightarrow_*F$ gives this implication at every
control: the two endpoint cases follow from the branch
implications. Existential elimination provides $M_a(A,B)$
and the implication gives $F$.

Given $A$ and $B$, selector endpoints and NC give $M_v(A,B)$
at an arbitrary control. Substitute $v=p(Sz)$ and use the
successor gate for conjunction introduction. For projections,
open its zero-gated clause and specialize at $1$ and $2$;
the two fixed computations give $B$ and $A$, respectively.

These rules prove existential numerical coverage:
\[
     (x=0)\vee_*\exists_*v\,(x=Sv).
\tag{Coverage}\label{eq:choice-coverage}
\]
Use NC. Zero takes the left injection of reflexivity, and a
successor takes the right injection with the evident witness.
Coverage gives equality cases even when the retained context
contains $x$; its elimination adds $x=0$ or $x=Sv$ rather
than silently changing that retained context.

Universal introduction follows from a fresh-parameter proof
by substituting zero and successors and forming the two
conjuncts. For elimination, take its $A(0)$ conjunct and
open the paired conjunct. Its base is a zero-gate theorem.
Coverage of the target chooses the zero instance from
$A(0)$ or a successor instance from the opened clause
and gate removal.

Finally, under $0=1$, computation and transport give
$s=C(0,s,t)=C(1,s,t)=t$ for arbitrary numerical terms.
Every equation follows, and external induction on formula
syntax with case DI gives every native formula. This
supplies falsity elimination. No excluded middle for
quantified formulas has entered the construction.

\subsection{Whole-domain realizations and both formula composites}

For a standard first-order reading, let $T_{\rm eq}$ have
intuitionistic equality logic, the same numerical signature
and computation equations, numerical coverage, separation
$(0=1)\to\bot$, and contextual equation induction.
Its interpretation of the native syntax is
\[
 (s=t)^\circ=(s=t),\qquad
 (\mathsf D_{x,y}B)^\circ
   =\exists x\bigl(B^\circ(x,0)\to\forall y B^\circ(x,y)\bigr),
\]
\[
 \langle A;y.B\rangle^\circ
       =A^\circ\to\forall y B^\circ.
\]
The domain is the entire numerical domain and interpreted
equality is actual equality.

Case DI is sound in this realization by numerical coverage.
Progress DI is sound by contextual equation induction under
the conditional base. Use, DE, transport, and computation
have their displayed ordinary meanings. Conversely, the
native macros derive every ordinary logical rule,
\eqref{eq:choice-ei} derives equation induction, and
\eqref{eq:choice-coverage} derives coverage. Separation
translates to the identity implication
$(0=1)\Rightarrow_*(0=1)$.

For the essential paired formula composite,
\[
 \mathsf D_{x,y}B
 \dashv\vdash
 \exists_*x\bigl(B(x,0)\Rightarrow_*\forall_*y\,B(x,y)\bigr).
\tag{Paired composite}\label{eq:choice-paired-composite}
\]
Forwards, DE opens its clause at a fresh $a$; under $B(a,0)$
that clause supplies $B(a,y)$ at a fresh arbitrary $y$.
Introduce the universal, implication, and existential
macros and close DE. Backwards, eliminate the existential
at a fresh $a$. Introduce the paired formula with witness
$a$ by case DI: the retained base, supplied implication,
and universal elimination at $Sz$ give its successor premise.

The ordinary interpretations of the six logical macros
have their ordinary round trips. For conjunction, its true
base and the instances $1,2$ recover both operands;
the converse uses selector cases. The universal has all
successor instances in its paired conjunct and zero in
its first conjunct; coverage gives exactly the ordinary
universal. Implication, disjunction, and existential
quantification use their selector or gate endpoints.
Falsity uses the explicit ordinary separation bridge.
External induction on finite formulas, together with
\eqref{eq:choice-paired-composite}, gives both formula
composites. Clauses correspond in both directions to
$A\Rightarrow_*\forall_*yB$ by the derived logical rules.

The numerical composites are literal identities with
internally provable graph $x=y$. No standard-numeral
restriction, retract, quotient, or equality reflection
assumption is being inserted at this stage. The separately
supplied arithmetic bridge identifies $T_{\rm eq}$ with
full iPRA; the present change preserves that bridge
unchanged because Theorem~\ref{thm:choice-equivalence}
preserves every native sequent.

\subsection{What the theorem does and does not claim}

The new inference interface retains one ordinary assumption
at every paired introduction. It preserves the entire
consequence relation, not just equations or closed
theorems. Consequently the established full-iPRA
interpretations and equiconsistency transfer unchanged.

This result is different from retaining only progress
introduction and strengthening closure elimination.
Here the original elimination rule is retained, and
the case and progress introductions coexist. It is also
different from allowing arbitrary compound progress,
which would supply stronger induction.

A complete-checker comparison must still count the
certificate choice, its equation-root restriction,
substitution, context handling, computation, and the
public input/output interface. The mathematical theorem
alone is not a byte-count or minimality theorem.
The inference-node bound is distinct from source-size
measurements and from serialized certificate size.
The proof is a written universal proof; finite regression
tests do not constitute a Lean/Rocq formalization of it.



\section{Arithmetic strength and both ordinary interpretation composites}
\label{sec:ordinary-bridge}

We spell out the arithmetic bridge to avoid identifying standard numerical
agreement with bi-interpretability. First elaborate every fixed iterator
into an ordinary function symbol. At stage zero the symbols are $0,S$.
Given a fixed earlier-stage term $h(u,\vec p)$, introduce
$i_h(n,a,\vec p)$ with the two equations
\[
i_h(0,a,\vec p)=a,\qquad
i_h(Sn,a,\vec p)=h(i_h(n,a,\vec p),\vec p).
\]
Only finite descriptions occur; the domain contains no code for a
universal evaluator. Formal parameters are distinct when a description
is declared, while actual arguments can coincide.

Let $T_I$ be ordinary intuitionistic first-order logic with actual
equality and genuine falsity, these staged equations, contextual
induction for literal equations, and
\[
\forall x\,(x=0\vee\exists y\,x=Sy),
\qquad (0=S0)\to\bot.
\]
The second sentence is the falsity bridge. It must be stated: positive
numerical equations by themselves admit a singleton structure. Native
explosion from $0=S0$ and a translation of ordinary falsity to that
equation do not authorize silently dropping the bridge from $T_I$.

The native-to-ordinary translation is
\[
(s=t)^\circ=(s=t),\qquad
(\mathsf D_{x,y}B)^\circ=
\exists x\bigl(B^\circ(x,0)\to\forall y B^\circ(x,y)\bigr).
\]
A clause $\langle A;y.B\rangle$ becomes $A^\circ\to\forall yB^\circ$.
For the auxiliary closure used by the compact checkers,
$(\mathsf Q_yB)^\circ=B^\circ(0)\to\forall yB^\circ$.
The interpretation domain is all elements, defined by $x=x$;
equality is actual equality.

Case DI is sound because numerical coverage combines its retained base
with its successor proof. Progress DI is sound by literal-equation
induction under a temporarily assumed base. Use, DE and transport are
ordinary logical inferences under the translations. Conversely the
derived logical operations, equation induction, cases and explosion in
Section~\ref{sec:choice-introduction} recover the rules and axioms of
$T_I$. An auxiliary closure is interderivable with
$\mathsf D_{v,y}B(y)$ for unused $v$: in one direction use case DI with
the closure and its base; in the other eliminate the paired formula
and return its supplied closure. Thus formulas and context clauses have
the stated meanings in both presentations.

\subsection{The arithmetic dependency order}

The scalar update has no recursion-counter argument. We do not conceal
one in an encoding. The inherited internal proofs, reproduced in the
appendices, establish the required arithmetic in this order:
\begin{enumerate}
\item Using equation induction and numerical cases, construct a fixed
scalar predecessor term and prove $P0=0$ and $P(Sn)=n$.
Existential coverage then gives $x=0\vee x=S(Px)$: in the successor
case $x=Sy$, apply $P$ to obtain $Px=y$ and substitute.
\item Construct truncated subtraction and a numerical characteristic
of actual equality. Prove its reflection on the entire original domain.
\item Compile each fixed quantifier-free formula into its characteristic
equation and recover every contextual quantifier-free induction instance.
\item With that induction available, construct pairing and projections
by a bounded square-root search and compile every fixed general
primitive-recursive description into a scalar iterator term.
\end{enumerate}
No pairing or faithful equality detector is assumed to obtain
predecessor. No appeal to standard-model correctness replaces an
internal equality. Every induction fixes its passive parameters;
diagonal substitution occurs only into a previously proved identity.

For a target recursion
\[
f(0,\vec p)=g(\vec p),\qquad
f(Sn,\vec p)=h(n,f(n,\vec p),\vec p),
\]
let $\pi,L,H$ be the proved pairing and projections. The compiler uses
\[
V(n,\vec p)=I\bigl(n,\pi(0,g^*(\vec p));
u.\pi(SL(u),h^*(L(u),H(u),\vec p))\bigr),
\qquad f^*(n,\vec p)=H(V(n,\vec p)).
\]
Equation induction proves $L(V(n,\vec p))=n$; the two projection
identities then prove the target defining equations. The construction
is external recursion on each finite PR description, not an internal
quantification over arbitrary functions.

\begin{theorem}[Whole-domain bi-interpretability and equiconsistency]
The ordinary realization $T_I$ and full first-order iPRA are
bi-interpretable. In both theories $x=y$ defines the isomorphism from
the original domain to its double interpretation. The native Choice-DI
calculus is therefore a faithful proof presentation of a theory
bi-interpretable and equiconsistent with iPRA.
\end{theorem}
\begin{proof}
Interpret an iterator in iPRA by primitive recursion that ignores its
counter argument in the update. iPRA proves its equations, literal
equation induction, coverage and separation. In the reverse direction
use the scalar compiler. The predecessor and equality-reflection proofs
supply the successor axioms and every quantifier-free induction instance.
All ordinary logical rules are preserved.

For two iterators with the same seed and provably equal updates,
equation induction on the equality of their values proves uniqueness at
every counter, with all parameters fixed. The corresponding argument
applies to ordinary primitive recursion. Induct externally on a finite
description: its double translation satisfies its original defining
equations, so internal uniqueness identifies the two functions. This
also proves coherence when substitution changes the staged syntactic
description of an update.

The domain and equality are unchanged; hence the identity graph is
internally total, bijective and preserves every interpreted operation.
The displayed ordinary meanings of the logical macros and the paired
binder give both formula composites by external induction on finite
formula syntax. Clauses are similarly equivalent to their conditional
universal meanings. The explicit falsity bridge preserves contradiction
in both directions. Effective translations of finite proofs give
equiconsistency, without a claim that either theory proves its own
consistency.
\end{proof}



% Adapted with attribution from the pinned unparenthesized one-hole archive.
\section{Internal scalar predecessor bootstrap}
\label{sec:hybrid-bootstrap}

This scalar bootstrap is inherited from the supplied scalar development~\cite{erasure} and is reproduced to keep the scalar checker self-contained. It is not a new contribution of this paper. The proof uses a factor map fixing odd states and a saturated-iteration construction of predecessor.

The available arithmetic terms are $0$, $S$, and
\[
 I(n,a;u.h),\qquad
 I(0,a;u.h)=a,\qquad
 I(Sn,a;u.h)=h[I(n,a;u.h)/u].
\]
The update binds only $u$.  All other variables in it are fixed numerical
parameters.  We use equality reasoning, the already reconstructed
zero/successor cases, and induction on a single literal equation.
Every induction below keeps its other parameters fixed.  Predecessor,
successor injectivity, pairing, and induction on arbitrary
quantifier-free formulas are not assumptions of the argument.

\subsection{Elementary terms and two finite cycles}

Write $1=S0$ and use larger numerals as successor abbreviations.  Define
\[
 C(n,a,b)=I(n,a;u.b),\qquad
 Z(n)=C(n,1,0),\qquad
 a+b=I(b,a;u.Su).
\]
The update variable of $C$ is absent from its body.  Consequently
\[
 C(0,a,b)=a,\qquad C(Sn,a,b)=b.
\]
Numerical cases show that $Z(n)$ is $0$ or $1$.
The usual commutative-monoid equations for addition follow by literal
equation induction.  For example, $0+b=b$ and $Sa+b=S(a+b)$ are
proved by induction on $b$; these yield commutativity by a further
equation induction, and associativity follows by induction on the last
summand.  We use these equations to reorder finite sums.  The notation
$2a$ means $a+a$ and $3a$ means $a+a+a$; no multiplication operation
is needed for this bootstrap.

Put
\[
 N(n)=I(n,0;u.Z(u)),\qquad
 j(u)=Z(u)+2N(u),\qquad
 M(n)=I(n,0;u.j(u)).
\]
The range of $N$ is contained in $\{0,1\}$: its zero value is $0$,
and its successor value is a value of $Z$.  This uses numerical cases
on its counter, not induction on a disjunction.  Hence
\begin{equation}
 N(n)+Z(N(n))=1,\qquad
 Z(Z(N(n)))=N(n),\qquad N(n+2)=N(n).
 \label{eq:hybrid-parity}
\end{equation}
Here the last equation follows from two successor computations and
the preceding equation.  Literal-equation inductions give
\begin{equation}
 N(2n)=0,\qquad N(2n+1)=1,\qquad N(3n)=N(n).
 \label{eq:hybrid-parity-scaling}
\end{equation}
For the first equation the successor step adds two to the argument of
$N$; the second follows by one computation.  In the third, increasing
$n$ adds three to the left argument and therefore applies $Z$ once to
both sides of the induction equation.

Finite computations now give
\[
 j(0)=1,\qquad j(1)=2,\qquad j(2)=0.
\]
The proof below needs just these three values and
$M(Sn)=j(M(n))$; it does not assume a division algorithm or a theorem
characterizing $M$ as remainder modulo three.

\subsection{A ramp with a global scaling identity}

Define
\[
 G(n)=I(n,0;u.u+M(u)+1).
\]
The central invariant is
\begin{equation}
 M(G(n))=N(n).
 \label{eq:hybrid-ramp-invariant}
\end{equation}
Its zero case is immediate.  For its successor step, retain the
single induction equation and split the Boolean value $N(n)$.
If it is $0$, the next $G$ state is $G(n)+1$; its $M$ value is
$j(0)=1$, equal to $N(Sn)$.  If it is $1$, the next $G$ state
is $G(n)+2$; two computations of $M$ change its value from $1$
to $2$ and then to $0$, again equal to $N(Sn)$.  This proves
the invariant by literal-equation induction.

It follows that
\[
 G(Sn)=G(n)+N(n)+1,\qquad G(n+2)=G(n)+3.
\]
Equation induction therefore yields
\begin{equation}
 G(2n)=3n.
 \label{eq:hybrid-ramp-double}
\end{equation}
A separate literal-equation induction proves
\begin{equation}
 G(3n)=3G(n)+N(n).
 \label{eq:hybrid-ramp-triple}
\end{equation}
For completeness, its successor calculation is
\begin{align*}
 G(3Sn)
 &=G(3n)+N(3n)+4
   =3G(n)+2N(n)+4,\\
 3G(Sn)+N(Sn)
 &=3G(n)+3N(n)+3+Z(N(n)).
\end{align*}
The expressions agree in both Boolean cases for $N(n)$.
The base is $0=0$.  In particular,
\eqref{eq:hybrid-ramp-triple} was proved internally; it was not
inferred from an intended floor-function interpretation of $G$.

Set
\[
 H(u)=C(N(u),G(u),u).
\]
The crucial choice is to retain $u$ when its parity is $1$.
The preceding equations give
\begin{equation}
 H(2u)=3u,\qquad H(1)=1,\qquad H(3u)=3H(u).
 \label{eq:hybrid-factor-map}
\end{equation}
To check the last identity, use $N(3u)=N(u)$ and split the Boolean
$N(u)$.  At $0$, both sides are $3G(u)$ by
\eqref{eq:hybrid-ramp-triple}; at $1$, both are $3u$.
This argument does not require that an arbitrary number has already
been represented as $2v$ or $2v+1$.

\subsection{A comparator from scalar iteration}

For every fixed update $h$, equation induction on $n$, keeping $m,a$
and all parameters fixed, proves the semigroup identity
\[
 I(m+n,a;h)=I(n,I(m,a;h);h).
\]
The base is the zero computation and the step applies $h$ to the
induction equation. This identity will also be used in the later
pairing construction.


For a fixed update $h$, equation induction on $y$, with its initializer
$a$ and all parameters of $h$ fixed, proves the rotation identity
\begin{equation}
 I(Sy,a;h)=I(y,h(a);h).
 \label{eq:hybrid-rotation}
\end{equation}
Both sides have base $h(a)$, and their successor equations apply the
same $h$ to the previous equality.  Here and below $I(y,a;h)$ is
an abbreviation for one fixed finite scalar term, not quantification
over functions.

By \eqref{eq:hybrid-factor-map}, another equation induction gives
\begin{equation}
 I(y,3a;H)=3I(y,a;H).
 \label{eq:hybrid-iterator-scaling}
\end{equation}
Its base is $3a=3a$.  In the step the fixed-parameter induction
equation rewrites the new left side as
\[
 H(3I(y,a;H))=3H(I(y,a;H)).
\]
The initializer parameter stays $a$ throughout this induction.
Equation induction also gives $I(y,1;H)=1$, since $H(1)=1$.

Define
\[
 E(x)=I(x,1;u.2u),\qquad
 L(x,y)=N\bigl(I(y,E(x);H)\bigr).
\]
Then the following entire recursion specification is provable:
\begin{equation}
 L(0,y)=1,\qquad L(Sx,0)=0,\qquad L(Sx,Sy)=L(x,y).
 \label{eq:hybrid-comparator}
\end{equation}
The first equation uses $E(0)=1$ and the fixed-point identity for
$H$.  The second uses $E(Sx)=2E(x)$ and $N(2a)=0$.
For the third,
\begin{align*}
 L(Sx,Sy)
 &=N\bigl(I(Sy,2E(x);H)\bigr)\\
 &=N\bigl(I(y,H(2E(x));H)\bigr)\\
 &=N\bigl(I(y,3E(x);H)\bigr)\\
 &=N\bigl(3I(y,E(x);H)\bigr)\\
 &=N\bigl(I(y,E(x);H)\bigr)=L(x,y).
\end{align*}
The successive steps use rotation, $H(2a)=3a$, iterator scaling,
and $N(3a)=N(a)$.  Thus $L$ has the comparison equations at all
arguments internally.  Its standard-model meaning as the indicator
of $x\le y$ is motivation, not a premise of the following argument.

\subsection{Saturated iteration gives predecessor}

We now need only the three equations
\eqref{eq:hybrid-comparator}.  Define
\[
 B_m(u)=C(L(Su,m),u,Su),\qquad
 W(t,m)=I(t,0;u.B_m(u)).
\]
The subscript in $B_m$ records a fixed parameter; $B_m$ is an open
term abbreviation, not a function variable.

Numerical cases on $e$ prove the selector equation
\[
 C(e,Sa,Sb)=S C(e,a,b).
\]
Equivalently one may use equation induction on $e$, whose induction
hypothesis is unused.  Therefore
\begin{align}
 B_{Sm}(Su)
 &=C(L(SSu,Sm),Su,SSu)\notag\\
 &=C(L(Su,m),Su,SSu)=S B_m(u),
 \label{eq:hybrid-conjugacy}\\
 B_{Sm}(0)&=C(L(S0,Sm),0,S0)=C(L(0,m),0,S0)=S0.
 \label{eq:hybrid-conjugacy-base}
\end{align}

Literal-equation induction on $t$, with $m$ fixed, proves
\begin{equation}
 W(St,Sm)=S W(t,m).
 \label{eq:hybrid-saturated-shift}
\end{equation}
The base is \eqref{eq:hybrid-conjugacy-base}.  The successor step is
\begin{align*}
 W(SSt,Sm)
 &=B_{Sm}(W(St,Sm))\\
 &=B_{Sm}(S W(t,m))\\
 &=S B_m(W(t,m))=S W(St,m).
\end{align*}
Only the second line uses the current induction hypothesis, at the
same fixed $m$.

Two subsequent equation inductions give
\begin{equation}
 L(Sn,n)=0,\qquad W(n,n)=n.
 \label{eq:hybrid-diagonals}
\end{equation}
The first uses the comparator recurrence.  The second uses
\eqref{eq:hybrid-saturated-shift} in its successor step.  The
substitution $t=m=n$ is made into that already proved universally
parameterized equation; no parameter of an active induction hypothesis
is changed.  Consequently
\begin{equation}
 W(Sn,n)=B_n(n)=C(L(Sn,n),n,Sn)=n.
 \label{eq:hybrid-extra-step}
\end{equation}

Finally define the finite scalar term
\begin{equation}
 P(m)=I\bigl(m,0;u.C(L(SSu,m),u,Su)\bigr).
 \label{eq:hybrid-predecessor-term}
\end{equation}
Its zero equation is immediate.  To justify its successor equation
without assuming a separate extensionality rule for iterator bodies,
put
\[
 V(t,m)=I\bigl(t,0;u.C(L(SSu,Sm),u,Su)\bigr).
\]
The comparator recurrence identifies this update pointwise with
$B_m(u)$.  Equation induction on $t$, with $m$ fixed, gives
$V(t,m)=W(t,m)$: their base values are $0$, and the successor step
uses the pointwise update equation and then the induction equation.
Now \eqref{eq:hybrid-extra-step} yields
\[
 \boxed{\qquad P(0)=0,\qquad P(Sm)=V(Sm,m)=W(Sm,m)=m.\qquad}
\]

\subsection{Dependency boundary}

All terms in this section are built from $0,S,I$ by finite substitution.
All inductions have a single equation as their invariant.  Numerical
cases and equality transport are the already reconstructed logical
resources.  Successor injectivity may now be obtained by applying $P$
to an assumed equation $Sx=Sy$.  Equality tests, reflection of
quantifier-free formulas, pairing, and the recovery of general
primitive recursion can then be developed in the established order.
They were not used in order to manufacture predecessor.

This argument is a written internal derivation in the stated theory.
The finite implementation checks do not constitute a proof-assistant
formalization of this entire metatheorem.



% Standalone section.  Requires only amsmath and amssymb in the parent document.
% No theorem environments, macros, or bibliography keys are assumed.
% All arithmetic claims in this section are proved here.

\section{From literal-equation induction to full intuitionistic PRA}
\label{sec:eic-completion}\label{app:arithmetic}

This section supplies an unconditional completion theorem. It is adapted
with attribution from the current forward-left draft~\cite{forward}. In the
scalar application used here, every function occurring in the proof is
already a scalar term: predecessor is supplied by
Appendix~\ref{sec:hybrid-bootstrap}, the selector is
$C(n,a,b)=I(n,a;u.b)$, addition is $x+y=I(y,x;u.Su)$, and
truncated subtraction is $d(x,y)=I(y,x;u.Pu)$. All other functions below
are finite compositions of these. Thus this proof does not presuppose
that arbitrary primitive recursion has already been compiled. The QFI
conclusion first applies to the scalar language; the general PR compiler
comes afterward.  Literal-equation
induction, together with the primitive-recursive defining equations,
predecessor coverage, and separation of $0$ and $1$, proves induction for every
quantifier-free formula.  In particular, the proof does not assume a faithful
numerical equality test, stability or decidability of equality, successor
uniqueness, or any commutation law for an iterator.  We construct the equality
test and prove its reflection property for arbitrary elements of the
numerical domain.

\subsection{The exact hypotheses}

Work in ordinary intuitionistic first-order logic with equality and genuine
falsity $\bot$.  The language has one numerical sort and a symbol for each
fixed, finite primitive-recursive description.  Distinct descriptions are
not identified merely because they compute the same standard function.
Let $\mathsf{PRE}$ contain the defining equations for zero, successor,
projections, composition, and primitive recursion.  In particular, choose
descriptions $P$ and $C$ with equations
\begin{align}
 P(0)&=0, & P(Sx)&=x, \label{eq:eic-predecessor}\\
 C(0,a,b)&=a, & C(Sx,a,b)&=b. \label{eq:eic-selector}
\end{align}
Write $1=S0$.  In addition to $\mathsf{PRE}$, assume
\begin{align}
 \mathsf{Cov}:&\qquad
       \forall x\,\bigl(x=0\lor x=S(Px)\bigr),
       \label{eq:eic-coverage}\\
 \mathsf{Sep}:&\qquad 0=1\longrightarrow\bot,
       \label{eq:eic-separation}
\end{align}
and induction for each \emph{literal equation} $E(n,\vec p)$:
\begin{equation}
 \forall\vec p\,
 \left[
   \left(E(0,\vec p)\land
       \forall n\bigl(E(n,\vec p)\longrightarrow E(Sn,\vec p)\bigr)
   \right)
   \longrightarrow\forall n\,E(n,\vec p)
 \right].
 \label{eq:eic-induction}
\end{equation}
Here the schematic expression $E$ is a single equality of numerical
terms; the displayed universal quantifiers range only over numbers.
Every description and every instance of the schema is finite external
syntax.

The axiom formulation \eqref{eq:eic-induction} and the usual contextual
induction rule are interchangeable.  For example, if
\[
 \Gamma\vdash E(0,\vec p),
 \qquad
 \Gamma,E(z,\vec p)\vdash E(Sz,\vec p),
\]
where $z$ is absent from $\Gamma$ and the passive parameters, discharge the
second premise's equation and universally introduce $z$.  The axiom then
gives $\Gamma\vdash\forall n E(n,\vec p)$, which can be instantiated at any
target term.  Conversely, assume the antecedent of the axiom and apply the
contextual rule with fresh variables.  This uses the ordinary intuitionistic
logical rules and permits arbitrary, including quantified, formulas in
$\Gamma$.

\subsection{Constructor separation and two elementary tests}

Define the sign description and the countdown description by
\begin{equation}
 \operatorname{sg}(u)=C(u,0,1),
 \qquad
 K(u)=C(u,1,Pu).
 \label{eq:eic-sign-countdown}
\end{equation}
Their defining equations imply
\begin{align}
 \operatorname{sg}(0)&=0,
 &\operatorname{sg}(Sx)&=1,\label{eq:eic-sign-computation}\\
 K(0)&=1,
 &K(Sx)&=x.\label{eq:eic-countdown-computation}
\end{align}

\paragraph{Constructor laws.}
Successor is injective: applying $P$ to $Sx=Sy$ gives $x=y$ by
\eqref{eq:eic-predecessor}.  Successor never equals zero.  Indeed,
$Sx=0$ gives $1=0$ by applying $\operatorname{sg}$ and using
\eqref{eq:eic-sign-computation}; symmetry and
\eqref{eq:eic-separation} give $\bot$.  Consequently $1\ne0$, where
$A\ne B$ abbreviates $A=B\to\bot$.  None of these arguments uses
induction.

\paragraph{The sign test.}
For every $u$,
\begin{align}
 &\operatorname{sg}(u)=0\lor\operatorname{sg}(u)=1,
       \label{eq:eic-sign-range}\\
 &\operatorname{sg}(u)=0\quad\longleftrightarrow\quad u=0,
       \label{eq:eic-sign-zero}\\
 &\operatorname{sg}(u)=1\quad\longleftrightarrow\quad u\ne0.
       \label{eq:eic-sign-one}
\end{align}
To prove the range assertion, split \eqref{eq:eic-coverage}.  The branch
$u=0$ gives sign $0$, and the branch $u=S(Pu)$ gives sign $1$.
The implication $u=0\to\operatorname{sg}(u)=0$ is computation.
For its converse, the successor branch of coverage would make
$\operatorname{sg}(u)=0$ imply $1=0$, a contradiction; in that branch
explosion gives $u=0$ as well.  Similarly, if
$\operatorname{sg}(u)=1$ and $u=0$, computation gives $0=1$.
Conversely, $u\ne0$ excludes the zero branch of coverage, and the
successor branch computes the sign as $1$.

Coverage and \eqref{eq:eic-selector} also give, for all $u,v$,
\begin{equation}
 C(u,v,v)=v.
 \label{eq:eic-selector-idempotence}
\end{equation}
Both coverage branches compute to $v$, so this identity again needs no
induction.

\paragraph{The countdown test and the notation $W$.}
For every $u$,
\begin{equation}
 K(u)=0\quad\longleftrightarrow\quad u=1.
 \label{eq:eic-countdown-zero}
\end{equation}
If $u=1=S0$, \eqref{eq:eic-countdown-computation} gives $K(u)=0$.
Conversely, assume $K(u)=0$ and split coverage.  If $u=0$, then
$K(u)=1$, contradicting separation, and explosion gives $u=1$.
If $u=S(Pu)$, computation gives $K(u)=Pu$; hence $Pu=0$ and
$u=S0=1$.

If a presentation uses a description named $W$ with defining equations
\[
 W(0)=1,\qquad W(Sx)=x,
\]
then it proves $K(u)=W(u)$ for every $u$ already from coverage.
In the zero branch both sides are $1$; in the branch $u=S(Pu)$ both
sides are $Pu$.  Thus $K$ and $W$ are provably equal on the whole domain,
even if they are distinct descriptions.  This equality does not appeal to
standard-model agreement or to uniqueness of an arbitrary recurrence.

\subsection{Four elementary uses of equation induction}

Choose descriptions for addition and truncated subtraction, with recursion
on the second argument:
\begin{align}
 x+0&=x, &x+Sy&=S(x+y),\label{eq:eic-addition}\\
 d(x,0)&=x, &d(x,Sy)&=P(d(x,y)).\label{eq:eic-subtraction}
\end{align}
Using these equations and \eqref{eq:eic-induction}, we now derive
\begin{align}
 0+y&=y,\label{eq:eic-left-zero}\\
 (Sx)+y&=S(x+y),\label{eq:eic-left-successor}\\
 d(Sx,Sy)&=d(x,y),\label{eq:eic-cancel-successors}\\
 d(Sa,a)&=1.\label{eq:eic-diagonal}
\end{align}
Each induction below has a single equation as its induction formula.

For \eqref{eq:eic-left-zero}, induct on $y$.  The base is $0+0=0$.
The step is $0+Sy=S(0+y)=Sy$.

For \eqref{eq:eic-left-successor}, hold $x$ fixed and induct on $y$.
At zero both sides are $Sx$.  At a successor the induction hypothesis
gives
\[
 (Sx)+Sy=S\bigl((Sx)+y\bigr)
       =S\bigl(S(x+y)\bigr)=S(x+Sy).
\]
Combining \eqref{eq:eic-left-successor} with the defining equation for
addition gives the useful rearrangement
\begin{equation}
 (Sy)+r=y+Sr.
 \label{eq:eic-shift-successor}
\end{equation}
Indeed, both sides are $S(y+r)$.

For \eqref{eq:eic-cancel-successors}, hold $x$ fixed and induct on $y$.
The base computes as
\[
 d(Sx,S0)=P(d(Sx,0))=P(Sx)=x=d(x,0).
\]
Applying $P$ to the induction equation gives the step:
\[
 d(Sx,S(Sy))=P(d(Sx,Sy))=P(d(x,y))=d(x,Sy).
\]

Finally, induct on $a$ for \eqref{eq:eic-diagonal}.  The base is
$d(S0,0)=S0=1$.  At a successor,
\[
 d(S(Sa),Sa)=d(Sa,a)=1,
\]
using \eqref{eq:eic-cancel-successors} and the induction hypothesis.

\subsection{An equation invariant that reflects actual equality}

The remaining arithmetic ingredient is the following equation, valid for
all $x,y$:
\begin{equation}
 C\bigl(\operatorname{sg}(d(x,y)),\ x,\ y+d(x,y)\bigr)=x.
 \label{eq:eic-residual-invariant}
\end{equation}
Its meaning is that a nonzero subtraction residual reconstructs the
original argument by addition.  The selector expresses this assertion as
one literal equation, so equation induction applies directly.

Hold $x$ fixed and induct on $y$.  At $y=0$, subtraction gives residual
$x$, and \eqref{eq:eic-left-zero} gives $0+x=x$.  The desired equation
is consequently $C(\operatorname{sg}(x),x,x)=x$, which is
\eqref{eq:eic-selector-idempotence}.

For the step, assume \eqref{eq:eic-residual-invariant} at $y$ and
abbreviate
\[
 r=d(x,y),\qquad r'=Pr=d(x,Sy).
\]
These are local term abbreviations, not additional function symbols.
Split coverage for $r'$.

If $r'=0$, then $\operatorname{sg}(r')=0$, and the successor instance
of \eqref{eq:eic-residual-invariant} reduces immediately to $x=x$.

In the other branch, $r'=S(Pr')$, so constructor separation gives
$r'\ne0$.  Split coverage for $r$.  Its zero branch would imply
$r'=P0=0$, a contradiction, from which the desired step equation follows.
In its successor branch,
\[
 r=S(Pr)=Sr',\qquad \operatorname{sg}(r)=1.
\]
The induction equation therefore computes to $y+r=x$.  Meanwhile,
\eqref{eq:eic-shift-successor} gives
\[
 (Sy)+r'=y+Sr'=y+r=x.
\]
Since $r'$ is a successor, $\operatorname{sg}(r')=1$, so this is
exactly the branch selected in the successor instance of
\eqref{eq:eic-residual-invariant}.  This completes the equation induction.
All case distinctions came from coverage; no induction on an implication
or a disjunction was used.

\paragraph{Reflection at residual $1$.}
We have, for arbitrary $a,b$,
\begin{equation}
 d(Sa,b)=1\quad\longleftrightarrow\quad a=b.
 \label{eq:eic-residual-reflection}
\end{equation}
For the forward direction, substitute $x=Sa$, $y=b$ into
\eqref{eq:eic-residual-invariant}.  If $d(Sa,b)=1$, the sign is $1$
and the selector returns its second branch, giving
\[
 b+1=Sa.
\]
The defining equation for addition gives $b+1=S(b+0)=Sb$.  Thus
$Sb=Sa$, and applying $P$ gives $b=a$, hence $a=b$.
Conversely, $a=b$ and \eqref{eq:eic-diagonal} give
$d(Sa,b)=d(Sa,a)=1$.

Define the equality characteristic by the fixed description
\begin{equation}
 e(a,b)=\operatorname{sg}\bigl(K(d(Sa,b))\bigr).
 \label{eq:eic-equality-characteristic}
\end{equation}
The sign, countdown, and residual reflection lemmas give
\begin{align}
 e(a,b)=0&\quad\longleftrightarrow\quad a=b,
       \label{eq:eic-equality-reflection}\\
 e(a,b)=0&\ \lor\ e(a,b)=1.
       \label{eq:eic-equality-range}
\end{align}
In detail, $e(a,b)=0$ is equivalent to $K(d(Sa,b))=0$ by
\eqref{eq:eic-sign-zero}, which is equivalent to $d(Sa,b)=1$ by
\eqref{eq:eic-countdown-zero}, and then to $a=b$ by
\eqref{eq:eic-residual-reflection}.  The range assertion is
\eqref{eq:eic-sign-range} applied to $K(d(Sa,b))$.

The reflection property concerns the logic's actual equality throughout
the domain.  It was not assumed as part of a numerical comparison symbol,
and its proof did not restrict $a$ or $b$ to standard numerals.
It also yields decidable equality.  In the $e(a,b)=0$ branch,
\eqref{eq:eic-equality-reflection} gives $a=b$.  In the $e(a,b)=1$
branch, an assumption $a=b$ would give $e(a,b)=0$ and therefore $0=1$,
contradicting separation.  Hence
\begin{equation}
 \forall a\,\forall b\,\bigl(a=b\lor(a=b\to\bot)\bigr)
 \label{eq:eic-equality-decidable}
\end{equation}
is a theorem, not an extra hypothesis.

\subsection{Compilation of all quantifier-free formulas}

For each fixed quantifier-free formula $A(\vec p)$, construct a term
$b_A(\vec p)$ by external recursion on its finite syntax.  Take equality,
$\bot$, $\land$, $\lor$, and $\to$ as the primitive formula constructors;
negation is implication into $\bot$.  Set
\begin{align}
 b_{s=t}&=e(s,t), & b_{\bot}&=1,\label{eq:eic-boolean-atoms}\\
 b_{A\land D}&=C(b_A,b_D,1),\label{eq:eic-boolean-and}\\
 b_{A\lor D}&=C(b_A,0,b_D),\label{eq:eic-boolean-or}\\
 b_{A\to D}&=C(b_A,b_D,0).\label{eq:eic-boolean-implies}
\end{align}
Every term here is a finite composition of fixed primitive-recursive
descriptions; if desired, projections and composition give it a single
description symbol.  No internal quantification over formulas or programs
is involved.

An external induction on the syntax of $A$ proves, within the arithmetic
theory,
\begin{equation}
 \bigl(b_A=0\lor b_A=1\bigr)
 \quad\text{and}\quad
 \bigl(b_A=0\longleftrightarrow A\bigr).
 \label{eq:eic-boolean-correctness}
\end{equation}
Here are the full connective checks.  In every case $b_A=1$ implies
$\neg A$: if $A$ held, the already established implication
$A\to b_A=0$ would contradict separation.

\paragraph{Atomic formulas and falsity.}
For an equation, both claims are
\eqref{eq:eic-equality-reflection}--\eqref{eq:eic-equality-range}.
For falsity the range assertion is $1=0\lor1=1$, with its right disjunct
proved by reflexivity.  Also $1=0\to\bot$ follows from separation and
symmetry, while $\bot\to1=0$ follows by explosion.  Thus $1=0$ is
equivalent to $\bot$, as required.

\paragraph{Conjunction.}
Split $b_A=0\lor b_A=1$.  In the zero branch, $A$ holds and
$b_{A\land D}=b_D$.  Its range follows from the range of $b_D$,
and its being zero is equivalent to $D$, hence to $A\land D$ under
the known $A$.  In the one branch, $\neg A$ holds and
$b_{A\land D}=1$.  Its range is immediate.  If it were zero,
separation would give $\bot$, hence $A\land D$; conversely an
assumption $A\land D$ contradicts $\neg A$, and explosion gives
$b_{A\land D}=0$.

\paragraph{Disjunction.}
In the branch $b_A=0$, the formula $A$ holds and
$b_{A\lor D}=0$.  The range is immediate, and both $b_{A\lor D}=0$
and $A\lor D$ hold, proving their equivalence.  In the branch
$b_A=1$, we have $\neg A$ and $b_{A\lor D}=b_D$.  Its range follows
from the range of $b_D$.  If its value is zero, then $D$ and hence
$A\lor D$.  Conversely, from $A\lor D$, the $A$ alternative
contradicts $\neg A$ and yields $b_D=0$ by explosion, while the
$D$ alternative yields $b_D=0$ directly.  This proves the required
equivalence in both alternatives.

\paragraph{Implication.}
In the branch $b_A=0$, the formula $A$ holds and
$b_{A\to D}=b_D$.  Its range follows from the range of $b_D$.
If its value is zero, then $D$, which gives $A\to D$.  Conversely,
$A\to D$ and the known $A$ give $D$, hence $b_D=0$.
In the branch $b_A=1$, we have $\neg A$ and
$b_{A\to D}=0$.  Its range is immediate.  To prove $A\to D$,
assume $A$; this contradicts $\neg A$, and explosion gives $D$.
Thus both $b_{A\to D}=0$ and $A\to D$ hold in this branch, and
their equivalence follows.  This argument uses the proved two-value
range of the finite characteristic.  It does not use excluded middle
for arbitrary formulas or change the surrounding intuitionistic logic.

This completes the external syntactic induction proving
\eqref{eq:eic-boolean-correctness} for every quantifier-free formula.

\subsection{Recovery of induction and the full theory}

Let $A(n,\vec p)$ be any fixed quantifier-free formula.  Suppose an
arbitrary retained context $\Gamma$ proves its base and its induction
step:
\[
 \Gamma\vdash A(0,\vec p),
 \qquad
 \Gamma,A(z,\vec p)\vdash A(Sz,\vec p),
\]
with the usual freshness condition on $z$.  By
\eqref{eq:eic-boolean-correctness}, the base gives
$\Gamma\vdash b_A(0,\vec p)=0$.  Under the additional assumption
$b_A(z,\vec p)=0$, the same equivalence gives $A(z,\vec p)$,
the given step gives $A(Sz,\vec p)$, and the equivalence gives
$b_A(Sz,\vec p)=0$.  Apply literal-equation induction to
$b_A(n,\vec p)=0$ and then apply its reflection implication.
The result is $\Gamma\vdash A(t,\vec p)$ for every target term $t$.
Thus the theory proves the full contextual quantifier-free induction
rule, including in arbitrary quantified contexts.

\paragraph{Completion theorem.}
In the common language of fixed primitive-recursive descriptions, the
theory
\[
 \mathsf{PRE}+\mathsf{Cov}+\mathsf{Sep}+\mathsf{EI}
\]
has exactly the same sequents as ordinary first-order intuitionistic
primitive recursive arithmetic with all quantifier-free induction
instances.  The forward simulation follows from the recovered
quantifier-free induction, successor injectivity, and successor separation;
all primitive-recursive defining equations and all logical rules are
shared.  For the reverse simulation, literal-equation induction is an
instance of quantifier-free induction, and successor separation gives
$\mathsf{Sep}$.  Quantifier-free induction proves $\mathsf{Cov}$:
at zero choose its left disjunct, and at a successor use
$Sx=S(P(Sx))$, which follows from $P(Sx)=x$, for its right disjunct.

Both interpretations therefore use the full domain, defined by $x=x$,
and preserve every function symbol, equality, logical connective, and
quantifier.  Both composites have the identity isomorphism, with graph
$x=y$.  The theories are consequently identity-domain bi-interpretable
and equiconsistent.  This completion theorem is unconditional once the
literal-equation induction rule has been obtained; it has no dependency
on a classical successor-uniqueness reduction.



\section{Pairing, primitive recursion, and internal numerical composites}
\label{sec:iterator-coding}
This appendix is adapted with attribution from the progress-closure
arithmetic development~\cite{progress}. Its prerequisites have now been
proved: scalar predecessor, arithmetic reflection, and full quantifier-free
induction. Pairing is a derived proof tool; it is not part of any checker.

\subsection{Pairing from a bounded square-root search}

Once QFI is available, a familiar square-based injection gives a
particularly direct compiler. Write $a\le b$ for $d(a,b)=0$.
The usual linear-order and semiring identities follow by fixed
quantifier-free inductions from the successor, addition, and monus
equations. Multiplication itself is the scalar term
\[
                        a\cdot b=I(b,0;u.u+a).
\]
Only primitive-recursive arithmetic is involved in the following
bounded root search; there is no minimization over an unbounded domain.

For a fixed parameter $n$, put
\[
 U_n(v)=C(d((Sv)^2,n),Sv,v),\qquad
 R_n(i)=I(i,0;v.U_n(v)),\qquad q(n)=R_n(n).
\]
The update increments $v$ exactly when $(v+1)^2\le n$, and
otherwise retains it. QFI in $i$, with $n$ fixed, proves
\[
 R_n(i)\le i,\qquad R_n(i)^2\le n,\qquad
 R_n(i)=i\ \vee\ n<(S(R_n(i)))^2.              \tag{Root invariant}
\]
The base has $R_n(0)=0$. If the update increments, its test
contradicts the strict-upper-bound alternative in the third invariant.
Thus the old state equals $i$, and the new state equals $Si$;
its square bound is the increment test. If the update does not increment, the
strict upper square bound holds and the state remains fixed.
An already established strict upper square bound forces every
subsequent update to remain fixed. These are exactly the cases
needed for the three displayed invariants.

At $i=n$, either the strict upper square bound already holds or
$R_n(n)=n$, in which case $n<(n+1)^2$ supplies it. Thus
\[
                        q(n)^2\le n<(q(n)+1)^2.           \tag{Root}
\]
Monotonicity of squares shows uniqueness of a number satisfying
Root. All formulas just used are quantifier free in the scalar
language, with fixed $n$ during the induction.

Now define
\[
 \pi(a,b)=(a+b)^2+a,\qquad
 L(n)=d(n,q(n)^2),\qquad H(n)=d(q(n),L(n)).
\]
For $k=a+b$, one has $k^2\le k^2+a<(k+1)^2$ since $a\le k$.
Root uniqueness gives $q(\pi(a,b))=a+b$, and therefore
\[
                    L(\pi(a,b))=a,\qquad H(\pi(a,b))=b.
\]
Surjectivity of $\pi$ is unnecessary. The original numerical domain
is never replaced by the range of this coding; pairing is used only
to store the finite state of a fixed primitive recursion. This
standard square pairing is a convenient reconstruction here, not
a claimed new arithmetic basis.

\subsection{Compiling each primitive recursion}

Suppose the fixed target description is
\[
 f(0,\vec p)=g(\vec p),\qquad
 f(Sn,\vec p)=h(n,f(n,\vec p),\vec p).
\]
Assume by external structural induction that the earlier $g,h$
already have scalar translations $g^*,h^*$. Define
\[
 W(u,\vec p)=\pi(SL(u),h^*(L(u),H(u),\vec p)),
\]
\[
 V(n,\vec p)=I(n,\pi(0,g^*(\vec p));u.W(u,\vec p)),
 \qquad f^*(n,\vec p)=H(V(n,\vec p)).
\]
The update has only its state variable bound, with every $p_i$
fixed. Equation induction proves $L(V(n,\vec p))=n$.
The two projection equations and scalar computation then prove
both defining equations for $f^*$ internally. Initial functions
and composition translate directly. This completes the translation
of every fixed finite PR description without adding a universal
evaluator.

\subsection{Why both numerical round trips are internal}

Two fixed scalar iterators with the same seed and provably equal
updates are internally equal at all counters: use literal-equation
induction comparing their values. Their base agrees, and their
successor equations plus congruence prove the step. The same
uniqueness argument works for ordinary primitive recursion.

For an ordinary PR description, induct externally on its finite
construction. The double translation satisfies its original
defining equations, so internal equation induction identifies it
with the original operation. For a native scalar description,
earlier translated updates agree by the external induction
hypothesis; scalar uniqueness identifies the double-translated
iterator. These proofs hold at arbitrary internal counters.

Thus the domain relation $x=y$ is provably a total bijection
preserving every fixed numerical operation in both composites.
The formula constructions and round trips in
Section~\ref{sec:choice-introduction} give the corresponding formula
composites. Both interpretations preserve falsity, taking native
$0=1$ to an ordinary contradiction and ordinary $\bot$ to native
$0=1$. Finite proof translations therefore establish equiconsistency.



\clearpage
\section{Complete normalized acceptance checker}
\label{app:checker}

The following is the complete measured module. Documentation and comments
are removed by the measurement convention; every executable definition
and the command-line entry point are present. The readable source and
the optional bounded version are separate files in the archive.

\begin{lstlisting}
import json
import sys

def need(condition):
    if not condition:
        raise ValueError('invalid certificate')

def walk(a, f, d=0, form=None):
    if type(a) is int:
        need(not form)
        return f(a, d)
    need(type(a) is list)
    size = len(a)
    need(form is None or (size in ((1, 2) if form else (2, 3)) and (form or size == 3 or a[0] == 0)))
    a = [walk(b, f, d + (2 if size == 1 else int(size == 3 and i == 2)), size == 1) for i, b in enumerate(a)]
    if size == 3:
        c, b, h = a
        if c == 0:
            return b
        if type(c) is list and len(c) == 2:
            return sub(h, walk([c[1], b, h], lambda i, d: i))
    return a

def syntax(a, n, form=False):

    def variable(i, d):
        need(-n <= i <= d)
        return i
    need(walk(a, variable, form=form) == a)

def shift(a, k):
    return walk(a, lambda i, d: i + k if i > d else i)

def sub(a, t, depth=0):
    return walk(a, lambda i, d: shift(t, d) if i == d + 1 else i - (i > d + 1), depth)

def derive(packet, g):
    n = len(g)
    code, arg, branch = packet
    if type(code) is int:
        need(0 <= code < len(g))
        syntax(arg, n - 1)
        base, body = g[code]
        if base:
            need(derive(branch, g) == base)
        a = sub(body, arg)
    elif code:
        syntax(arg, n - 1)
        body = sub(code, arg, 1)
        hyp = sub(body, -n if len(body) == 2 and arg == 0 else 0)
        need(derive(branch, g + [([], hyp)]) == sub(body, [0, -n]))
        a = [code]
    else:
        need(code == [])
        body, = derive(arg, g)
        body = sub(body, -n, 1)
        base = sub(body, 0)
        row = sub(body, [0, 0])
        a = walk(derive(branch, g + [(base, body)]), lambda i, d: row[1] if base == [0, 0] and row[0] == i == -n else i)
    syntax(a, n - 1, True)
    return a

def check(data):
    need(type(data) is dict and set(data) == {'goal', 'proof'})
    syntax(data['goal'], 0, True)
    need(derive(data['proof'], [([], [1, 1])]) == data['goal'])
if __name__ == '__main__':
    try:
        check(json.load(sys.stdin))
    except Exception:
        print('REJECT')
        sys.exit(1)
    print('ACCEPT')
\end{lstlisting}

\section{Reproduction and file map}

With Python 3.12.14, run the following from the extracted directory:
\begin{lstlisting}[language={}]
python3 reproduce.py
python3 code/checker.py < tests/positive/successor_congruence.json
python3 code/checker.py < tests/negative/goal_false.json
python3 build_pdf.py
\end{lstlisting}
The first selected example prints \texttt{ACCEPT}; the second prints
\texttt{REJECT}. The reproduction entry point verifies the delivered
manifest, replays positive and negative certificates and proof translations,
checks the bounded variants, and recomputes complete-module measurements.
The PDF build uses an installed pdfLaTeX distribution with shell escape
disabled. A freshly compiled PDF may differ in metadata and therefore in
its file hash; the delivered manifest describes the original payload.

\begin{longtable}{@{}p{.31\textwidth}p{.65\textwidth}@{}}
\toprule
File or directory & Contents \\
\midrule
\endhead
\path{ipra_first_instance.tex} & Standalone LaTeX source containing this entire report. \\
\path{ipra_first_instance.pdf} & Compiled report. \\
\path{paper/} & Editable source sections used to assemble the standalone source. \\
\path{code/} & Complete selected and matched control checkers, with and without fuel. \\
\path{tests/} & Positive, negative and source certificates; untrusted proof constructors, translators and independent replay. \\
\path{reference/} & Attributed foundation fragments and exact comparison checker snapshots. \\
\path{evidence/} & Measurements, normalized and gzip sources, verification logs and provenance. \\
\path{FORMAT.txt} & Public JSON grammar and scope conventions. \\
\path{CLAIMS.json} & Explicit positive claims and limitations. \\
\path{SOURCE_PROVENANCE.json} & Pinned repository and mathematical source identities. \\
\path{MANIFEST.json} & SHA-256 checksums of every other payload file. \\
\bottomrule
\end{longtable}

\begin{thebibliography}{9}

\bibitem{baseline}
Vladimir Reshetnikov, \emph{ProveIt: Concise presentations of intuitionistic
primitive recursive arithmetic}, snapshot
\texttt{17898ea7e6a1300dfdcad71040d5daeed67916d4},
\path{Logic/PrimitiveRecursiveArithmetic/ConciseIPRA}.
\url{https://github.com/VladimirReshetnikov/ProveIt/tree/17898ea7e6a1300dfdcad71040d5daeed67916d4/Logic/PrimitiveRecursiveArithmetic/ConciseIPRA}.

\bibitem{witness}
\emph{iPRA with equality at witness discharge},
\path{docs/incoming/ipra_witness_discharge_2026-11-11.zip}
in~\cite{baseline}, especially the complete native specification,
naming-matrix proof, zero-witness repair, and normalization appendix.
The archive's own original snapshot is
\texttt{0ebd0c6b358f38bb0ad2883faf775421599c8aae}.

\bibitem{normalization}
\emph{Full normalization}, \path{Research/full_normalization}
in~\cite{baseline}; normalization results as reproduced and attributed
in~\cite{witness}. The present bridge additionally spells out a
termination justification for strict normalization of an unused
predecessor iterate.

\bibitem{uniformbinding}
\emph{Uniform binding for intuitionistic primitive recursive arithmetic},
\path{docs/incoming/ipra_uniform_binding_2026-11-11.zip}
in~\cite{baseline}. Exact checker source retained for comparison.

\bibitem{erasure}
\emph{Internal scalar predecessor bootstrap}, the attributed development
in \path{Research/certified_zero/reference/scalar_predecessor.tex}
of~\cite{baseline}; reproduced here without a claim of new authorship.

\bibitem{forward}
\emph{From literal-equation induction to full intuitionistic PRA},
the attributed forward-left development in
\path{Research/certified_zero/reference/equality_reflection.tex}
of~\cite{baseline}; reproduced here.

\bibitem{progress}
\emph{Pairing, primitive recursion, and internal numerical composites},
the attributed scalar arithmetic development in
\path{Research/certified_zero/reference/primitive_recursion.tex}
of~\cite{baseline}; reproduced here.

\bibitem{khaniki}
Erfan Khaniki, \emph{Not all Kripke models of HA are locally PA},
\emph{Advances in Mathematics} 397 (2022), 108126.
Author manuscript dated 9 January 2021, Section 2.1.
\url{https://iuuk.mff.cuni.cz/~koucky/EPAC/papers/HA_Final.pdf}.
Used for the full first-order iPRA convention, not as a source of the
new inference restriction.

\end{thebibliography}
\end{document}
